%File: anonymous-submission-latex-2024.tex
\documentclass[letterpaper]{article} % DO NOT CHANGE THIS
\usepackage[submission]{aaai24}  % DO NOT CHANGE THIS
\usepackage{times}  % DO NOT CHANGE THIS
\usepackage{helvet}  % DO NOT CHANGE THIS
\usepackage{courier}  % DO NOT CHANGE THIS
\usepackage{amsmath}
\usepackage[hyphens]{url}  % DO NOT CHANGE THIS
\usepackage{graphicx} % DO NOT CHANGE THIS
\usepackage{kotex}
\urlstyle{rm} % DO NOT CHANGE THIS
\def\UrlFont{\rm}  % DO NOT CHANGE THIS
\usepackage{natbib}  % DO NOT CHANGE THIS AND DO NOT ADD ANY OPTIONS TO IT
\usepackage{caption} % DO NOT CHANGE THIS AND DO NOT ADD ANY OPTIONS TO IT
\frenchspacing  % DO NOT CHANGE THIS
\setlength{\pdfpagewidth}{8.5in} % DO NOT CHANGE THIS
\setlength{\pdfpageheight}{11in} % DO NOT CHANGE THIS
\usepackage{xcolor} % Required for text color
%\usepackage{graphicx}
\usepackage{algcompatible}
\usepackage{algpseudocode}
\usepackage{amssymb}
\newcommand{\nj}[1]{\textcolor{red}{#1}}
\newcommand{\hh}[1]{\textcolor{green}{#1}}
\newcommand{\jh}[1]{\textcolor{blue}{#1}}
\newcommand{\yr}[1]{\textcolor{magenta}{#1}}
\usepackage{adjustbox}
\newcommand{\etal}{\textit{et al}., }
\usepackage{adjustbox}
\newcommand{\ie}{\textit{i}.\textit{e}., }
\newcommand{\eg}{\textit{e}.\textit{g}.\ }
\usepackage{lipsum}
%
% These are recommended to typeset algorithms but not required. See the subsubsection on algorithms. Remove them if you don't have algorithms in your paper.
\usepackage{algorithm}
%\usepackage{algorithmic}
\usepackage{layouts}
%
% These are are recommended to typeset listings but not required. See the subsubsection on listing. Remove this block if you don't have listings in your paper.
\usepackage{newfloat}
\usepackage{listings}
\DeclareCaptionStyle{ruled}{labelfont=normalfont,labelsep=colon,strut=off} % DO NOT CHANGE THIS
\lstset{%
	basicstyle={\footnotesize\ttfamily},% footnotesize acceptable for monospace
	numbers=left,numberstyle=\footnotesize,xleftmargin=2em, % show line numbers, remove this entire line if you don't want the numbers.
	aboveskip=0pt,belowskip=0pt,%
	showstringspaces=false,tabsize=2,breaklines=true}
\floatstyle{ruled}
\newfloat{listing}{tb}{lst}{}
\floatname{listing}{Listing}
%
% Keep the \pdfinfo as shown here. There's no need
% for you to add the /Title and /Author tags.
\pdfinfo{
/TemplateVersion (2024.1)
}

% DISALLOWED PACKAGES
% \usepackage{authblk} -- This package is specifically forbidden
% \usepackage{balance} -- This package is specifically forbidden
% \usepackage{color (if used in text)
% \usepackage{CJK} -- This package is specifically forbidden
% \usepackage{float} -- This package is specifically forbidden
% \usepackage{flushend} -- This package is specifically forbidden
% \usepackage{fontenc} -- This package is specifically forbidden
% \usepackage{fullpage} -- This package is specifically forbidden
% \usepackage{geometry} -- This package is specifically forbidden
% \usepackage{grffile} -- This package is specifically forbidden
% \usepackage{hyperref} -- This package is specifically forbidden
% \usepackage{navigator} -- This package is specifically forbidden
% (or any other package that embeds links such as navigator or hyperref)
% \indentfirst} -- This package is specifically forbidden
% \layout} -- This package is specifically forbidden
% \multicol} -- This package is specifically forbidden
% \nameref} -- This package is specifically forbidden
% \usepackage{savetrees} -- This package is specifically forbidden
% \usepackage{setspace} -- This package is specifically forbidden
% \usepackage{stfloats} -- This package is specifically forbidden
% \usepackage{tabu} -- This package is specifically forbidden
% \usepackage{titlesec} -- This package is specifically forbidden
% \usepackage{tocbibind} -- This package is specifically forbidden
% \usepackage{ulem} -- This package is specifically forbidden
% \usepackage{wrapfig} -- This package is specifically forbidden
% DISALLOWED COMMANDS
% \nocopyright -- Your paper will not be published if you use this command
% \addtolength -- This command may not be used
% \balance -- This command may not be used
% \baselinestretch -- Your paper will not be published if you use this command
% \clearpage -- No page breaks of any kind may be used for the final version of your paper
% \columnsep -- This command may not be used
% \newpage -- No page breaks of any kind may be used for the final version of your paper
% \pagebreak -- No page breaks of any kind may be used for the final version of your paperr
% \pagestyle -- This command may not be used
% \tiny -- This is not an acceptable font size.
% \vspace{- -- No negative value may be used in proximity of a caption, figure, table, section, subsection, subsubsection, or reference
% \vskip{- -- No negative value may be used to alter spacing above or below a caption, figure, table, section, subsection, subsubsection, or reference

\setcounter{secnumdepth}{0} %May be changed to 1 or 2 if section numbers are desired.

% The file aaai24.sty is the style file for AAAI Press
% proceedings, working notes, and technical reports.
%

% Title

% Your title must be in mixed case, not sentence case.
% That means all verbs (including short verbs like be, is, using,and go),
% nouns, adverbs, adjectives should be capitalized, including both words in hyphenated terms, while
% articles, conjunctions, and prepositions are lower case unless they
% directly follow a colon or long dash
\title{Any-Way Meta Learning}
\author{
    %Authors
    % All authors must be in the same font size and format.
    Written by AAAI Press Staff\textsuperscript{\rm 1}\thanks{With help from the AAAI Publications Committee.}\\
    AAAI Style Contributions by Pater Patel Schneider,
    Sunil Issar,\\
    J. Scott Penberthy,
    George Ferguson,
    Hans Guesgen,
    Francisco Cruz\equalcontrib,
    Marc Pujol-Gonzalez\equalcontrib
}
\affiliations{
    %Afiliations
    \textsuperscript{\rm 1}Association for the Advancement of Artificial Intelligence\\
    % If you have multiple authors and multiple affiliations
    % use superscripts in text and roman font to identify them.
    % For example,

    % Sunil Issar\textsuperscript{\rm 2},
    % J. Scott Penberthy\textsuperscript{\rm 3},
    % George Ferguson\textsuperscript{\rm 4},
    % Hans Guesgen\textsuperscript{\rm 5}
    % Note that the comma should be placed after the superscript

    1900 Embarcadero Road, Suite 101\\
    Palo Alto, California 94303-3310 USA\\
    % email address must be in roman text type, not monospace or sans serif
    proceedings-questions@aaai.org
%
% See more examples next
}

%Example, Single Author, $\rightarrow$> remove \iffalse,\fi and place them surrounding AAAI title to use it
\iffalse
\title{Any-Way Meta Learning through node equivalence}
\author {
    Author Name
}
\affiliations{
    Affiliation\\
    Affiliation Line 2\\
    name@example.com
}
\fi

\iffalse
%Example, Multiple Authors, $\rightarrow$> remove \iffalse,\fi and place them surrounding AAAI title to use it
\title{Any-Way Meta Learning through node equivalence}
\author {
    % Authors
    First Author Name\textsuperscript{\rm 1},
    Second Author Name\textsuperscript{\rm 2},
    Third Author Name\textsuperscript{\rm 1}
}
\affiliations {
    % Affiliations
    \textsuperscript{\rm 1}Affiliation 1\\
    \textsuperscript{\rm 2}Affiliation 2\\
    firstAuthor@affiliation1.com, secondAuthor@affilation2.com, thirdAuthor@affiliation1.com
}
\fi


% This is only needed to show inline citations in the guidelines document. You should not need it and can safely delete it.
\usepackage{bibentry}
% END REMOVE bibentry

\begin{document}

\maketitle

\begin{abstract}
\jh{Meta-learning has showcased prowess in rapid adaptability, yet its reliance on fixed cardinalities contrasts starkly with human learning's inherent flexibility. This paper challenges the existing paradigm \nj{of meta-learning} by harnessing `label equivalence', emerging from stochastic numeric label assignments during episodic task sampling. We introduce the ``any-way" learning paradigm, an innovative model training approach free from fixed cardinality constraints. Surprisingly, this model not only matches but often outperforms traditional fixed-way models in terms of performance, convergence speed, and stability, disrupting long-held beliefs about domain generalization. To address the semantic information gap resulting from label equivalence, we further propose a mechanism to infuse semantic class information \nj{to the model}, enhancing the model's comprehension and functionality. Experiments on renowned architectures like MAML and ProtoNet affirm our method's effectiveness, pushing the boundaries of what defines ``true" meta-learning and setting the stage for a more versatile and semantically-aware future in the field.}
\end{abstract}

\section{Introduction}
Meta-learning, often referred to as `learning to learn', is a training strategy designed to enable rapid adaptation to new tasks with only a few examples. Despite the impressive performance of deep learning models on extensive datasets, they still fall short of human abilities when it comes to learning from a small number of instances. Various approaches have been proposed to tackle this challenge \cite{labelequiv:metadefinition}.

There are two main meta-learning approaches. The first approach focuses on creating an effective feature extractor, like ProtoNet \cite{episodic:MBML:Protonet}, which can accurately cluster unseen instances from the same class to enable the model to generalize well to new data. The second approach aims to find a suitable initialization point for the model, facilitating rapid adaptation to unseen tasks \nj{\cite{episodic:GBML:MAML}}. 

To imbue the model with meta-learning ability, both methods involve episodes. Each episode consists of multiple tasks and tests the model's meta-learning ability in the regime of few-shot learning
(Fig. \ref{fig:task_sampling}). This continuous exposure to new tasks helps the model improve its learning ability from limited data, which simulates the conditions under which humans learn.

\jh{\nj{Nonetheless, in contrast to the human's remarkable flexibility in classifying varying numbers of classes, the current learning} models remain rigid, restricted to the fixed `way' they were trained on. As evident in Table~\ref{table:intuition}, when a model trained on \nj{10-way-5-shot tasks} from \nj{miniImageNet \cite{dataset:miniimagenet}} is tested on a 3-way-5-shot cars \cite{dataset:cars} setting, its performance plunges to a mere 41$\%$. \nj{Considering} that random sampling would yield at least 33$\%$, this reveals a glaring deficiency in \nj{meta-learning}'s adaptability across different `way' settings.}

\jh{This sharp decline in performance is more than just a testament to the model's inflexibility; it also suggests that different `way' episodes might hail from different distributions. This variation between distributions touches upon a deeper, underlying problem: domain generalization. Essentially, to achieve true `any-way' learning, our models must address the broader challenge of \nj{performing effectively even in} disparate, unseen domains.}

\jh{The core objective of this paper is to highlight and address this challenge, advocating for an approach where task cardinality agnosticity becomes a foundational element in meta-learning.}

\begin{figure}[t!]
\centering
\includegraphics[width=0.4\textwidth]{figure/greenisthe.pdf}
\caption{Task sampling procedure of our \textit{any-way} meta-learning. Meta-learning learns through episodic learning. It has a mother dataset with a large number of classes ($M$) where each label represents each \textit{semantic class}. % (e.g. `turtle', `octopus', and `deer')}.
Each episode samples $N < M$ classes from this mother dataset and a random \textit{numeric class} ($\in [N]$) is assigned to each selected class. Note that in our any-way setting, $N$ is different for different tasks. If $N$ is fixed along episodes, it becomes conventional sampling procedure. %are assigned to random task indices injective manner.} 
}
\label{fig:task_sampling}
\end{figure} 

To deal with this problem, we propose an \textit{‘any-way’} meta-learning method based on the concept of ‘label equivalence’, which emerges from the task sampling process, an essential element in few-shot learning including meta-learning.
The task sampling process has a unique characteristic: classes are randomly sampled and assigned to (typically numeric) class labels within each episode. These labels are not assigned based on the inherent semantics of the classes; rather, they are determined by the requirements of the specific tasks in each episode. As a result, the same class can be assigned different labels in different tasks and episodes, and the meaning of a label can change from one task to another. As a result, a numeric class label in a given task does not necessarily represent a specific class\footnote{In this paper, we differentiate the `semantic class' from the `numeric class label' as shown in Fig. \ref{fig:task_sampling}.}. This distinction introduces a phenomenon we term `label equivalence', which suggests that all output labels are functionally equivalent to one another.

\begin{table}[t]
\caption{\jh{Performance Comparison between different number of ways, a model is trained within miniImageNet. number in paraenthethis indicates number of shots.}}
\centering
\begin{tabular}{c|cc|cc}
\hline
\# of ways (train)    & \multicolumn{1}{c|}{3} & 10    & \multicolumn{1}{c|}{5} & 10    \\ \hline
\# of ways (test)     & \multicolumn{2}{c|}{3} & \multicolumn{2}{c}{5} \\ \hline \hline
miniImageNet (5) & 76.27      & 62.45     & 64.92     & 55.2      \\ \hline
Cars (5)        & 59.95      & 41.47     & 47.73     & 32.38    \\ \hline
\end{tabular}
\label{table:intuition}
\end{table}

Examining label equivalence naturally brings up an important question: Given the equivalence of all labels or output nodes,
%If all labels or output nodes are equivalent, 
is it necessary 
to match the number of output nodes to the number of classes in each task? For instance, if a new task consists of three classes, we could select 3 nodes from a pool of, let's say,  
% among, for example, 
30 output nodes and designate them as numeric labels, optimizing solely for this selected path. This implies that there's no need to rigidly stick to a fixed number of numeric labels, potentially allowing us to further generalize the few-shot learning process. It is from this juncture that our discourse begins.

The primary contribution of this paper is the world-premiere proposition of 
%\jh{generalized} 
`any-way' few-shot learning, moving beyond the conventional fixed-way few-shot learning. Not only do we extend meta-learning to any-way, but we also provide an efficient, general learning algorithm that achieves performance comparable to or even better % or potentially even better
than fixed-way learning. Moreover, label equivalence further suggests that individual outputs cannot represent the semantic of a class. This arises from the fact that each few-shot task %, whether episode-dependent or not, 
contains no semantic information about its class, nor does it recognize or account for the classes it comprises. From this perspective, we can perceive supervised learning as solving a given task within an absolute coordinate system (one axis for one semantic class), while few-shot learning can be viewed as problem-solving within a relative coordinate system, where the uniqueness of each class is the only clue for learning. %the only known variable is the uniqueness of each \nj{class}.} 

Our second contribution deals with the challenges arising from the lack of semantic class information. In this paper, we introduce an algorithm that compensates for this issue by integrating semantic class information into the learning process. We show that our algorithm not only improves performance in environments where the distribution of the test set differs from that of the training set, but also enables the application of data-augmentation techniques used in supervised learning to few-shot learning tasks. We demonstrate the effectiveness of this approach in MAML and ProtoNet, two of the most representative methods of episode-based meta-learning.

% Lastly, we show that these techniques are also applicable to conventional episode-independent few-shot methods \cite{}. Although episode-independent methods can sometimes outperform episode-based learning, their inherent approach of addressing the problem through the data scarcity problem makes it difficult to regard them as true meta-learning techniques. \jh{This paper, \yr{however,} shows that we can even increase the performance of episode-independent learning by injecting meta-episode semantic information only from few episodes. In essence, %we can have our cake and eat it too: 
% we can harness the benefits of the meta-learning methodology for episode-independent methods, albeit with the trade-off of using a few known episodes.}
% %\nj{*** 마지막 문장의 의미가 이해가 안 됨.***}

% \jh{The thing is both shares same episode sampling procedure. As seen in Fig.~\ref{fig:tasksampling} They do not have any information of given episode(Task? Episode? 통일 필요). Also, each class can be assigned to different label in different episode. In this paper, we focus on this phenomenon. We will call this phenomenon as 'label equivalence'. Which means no label is more different than others. Every label is equivalent to others. As this phenomenon is more important to episode-dependent meta learning. We will focus more on this. However, we also deal with episode-independent meta learning by using episode-dependent feature to it.}

% \jh{Also, this task set sampling methodology lacks class information. Even if the task-label set is fixed, each label is only a match to the task-label in that task (episode), and the neuron does not store any special information. Let's compare this phenomenon to supervised learning.  This does not happen in supervised learning. If we assign class A to output 0, then output 0 represents class A. What this means is that the initial state of the output neurons before learning a new task is not different from each other. We will refer to this phenomenon as label equivalence.}
% \jh{In this paper, we focus on label-equivalence because it is a unique phenomenon in few-shot learning. In this paper, we analyze this phenomenon and propose a way to take advantage of it. As a result, our paper aims to provide one insight into few-shot learning based on the keyword label equivalence.}
% \jh{이를 살펴볼 때, 먼저 이전에 episode-learning based 모델의 경우 진짜 meta-learning ability를 가지고 있는 것에 가깝다고 했지만, task-label의 개수가 fixed되어 있다는 조건 때문에 진정한 meta-learning 능력을 가지고 있다고는 의문이 있음. 이는 애초에 desired way의 크기를 키운 다음에 더 작은 incoming task set에 대해서 inference하면 된느거 아닌가? 하는 의문이 생길 수 있지만, 이는 후술하겠지만 일반적으로 episoded-based 모델의 경우 학습한 way set 에 대해서만 성능이 나오고, 이와 다를 때 성능이 급격히 떨어지는 경향을 보임. }
% \jh{Task set을 sampling 하는 점의 특징은 이 뿐만이 아님, 왜냐하면 이렇게 task-label set을 고정해놓는다고 하더라도 , 각각의 label은 해당 task(episode)에서의 task-label을 맞추는 것일 뿐, 해당 neuron이 특별한 정보를 저장하는 것이 아니기 때문. 이 현상이 무엇인지 supervised learning과 비교해 보겠음.  supervised learning에서는 이러한 현상이 일어나지 않음. 만약 0번째 output에 class A를 assign해주었다면 0번째 output은 class A를 represent함. 이는 무엇을 의미하냐면, 새로운 task에 대해서 학습하기 이전의 초기 상태는 output neuron 서로 다르지 않다는 것임. 본 논문에서는 이러한 현상으f label equivalence로 지칭하겠음.}
% \jh{본 논문에서는 label-equivalence에 집중할 것인데, 왜냐하면 이러한 현상은 supervised learning의 일종이라고 볼 수 있는 few-shot learning의 고유한 현상임에도 불구하고, 이 equivalence 자체에 대해서 집중한 페이퍼는 거의 없었기 때문임. 본 논문에서는 이러한 현상을 분석하고, 또 이 현상에서 advantage를 얻을 수 있는 방법을 제시하고자 함. 결과적으로 우리의 논문은 few-shot learning에 대해서 label equivalence라는 키워드를 바탕으로 하나으ㅣ insight를 제공하고자 ㅎㅏa.}
% \jh{label equivalence를 살펴보았을 때, 다음과 같은 질문은 자연스러움. 어차피 각 neuron이 equivalent하다면, 굳이 매 task에서 neuron의 개수와 task label의 개수를 맞출 필요가 없지 않은가? 만약 새로운 task의 label set이 3개라면, 우리는 output neuron중에서 3개를 골라서 이를 task-label이라고 하고 이 selection path에 대해서만 optimization을 하면 됨. 그렇다면 우리는 굳이 fixed-number of task label에 집착할 필요가 없으며, 더 나아가 few-shot learning의 과정을 좀더 일반화할 수 있음. 우리의 이야기는 이곳에서 시작됨. 우리의 첫 번째 contribution은 기존의 fixed-way few shot learning이 아닌, any-way few shot learning을 제시한다는 것. 우리는 단순히 any-way로 늘리는 것 뿐 아니라, 매 way의 성능이 fixed-way의 성능과 유사한 수준으로 만들어주는 이에 맞는 적절하고, 일반화된 학습 알고리즘을 제시하며, 오히려 fixed-way보다 성능이 더 좋을 수도 있다는 것을 보여줌 (이부분은 아직 데이터가 다 안나와서.. 추후 맞춰서 수저ㅇ)}
% \jh{label equivalence는 또한 각각의 ㅐutput들이 class를 represent 할 수 없게 만듬. 이는 왜냐하면 episode-dependent하든, 아니든 간에 매 task 에 대한 정보는 전혀 없고, 또한 그 task가 실제로 어떤 class를 가지고 있는지는 전혀 상관이 없고, 또 모르기 때문. 이러한 현상을 놓고 본다면 supervised learning에 비해서 few-show classification task는 제공받는 정보의 차원이 하나 더 낮다고 생각할 수 있음. 비유하자면 supervised learning은 절대 좌표계 속에서 주어진 task를 푸는 거라면 few-shot learning의 경우는 각 task들이 서로 다르다는 사실만 알 고 있는, 상대 좌표계 속에서 문제를 푸는 것으로 보임. 본 논문에서는 해당하는 문제를 보완하여 학습 과정 속 ㅊlass-label information을 inject하는 알고리즘을 제시하였으며, 이 알고리즘이 training set과 distribution이 비슷한 test set이 아닌, distribution이 다른 더 practical한 환경에서도 마찬가지로 성능 개선을 늘릴 수 있음을 보여줌. 또한, 이를 통해서 few-shot learning task에 supervised learning에서 사용하는 기법ㅇ들을 적용할 수 있음을 보여줌. 위의 두 방법들에 대해서는 episode-based learning에서 가장 대표격인 MAML, ProtoNET에서 이것이 작동함을 보였음.}
% \jh{본 논문에서는 마지막으로 이러한 기법들이 episode-independent한 기법에 대해서도 적용 가능하다는 것을 보여줌. 기존의 episode-independent한 기법은 성능 면에서는 episode-based learning보다 도 좋을 수 있지만, 결국 이를 ㅇataset scarcity problem을 통해 문제를 풀기 때문에 이를 true meta learning 기법이라고 보기에는 무리가 있음. 본 논문에서는 특정한 제약조건을 추가하였을 때 모델이 previous episode에서 학습하여 크게 성능을 개선할 수 있음을 보여줌.}
% \jh{정리하자면, 본 논문에서는 task sampling method에 주목하였고, 이를 통하여 meta learning 기법들을 더욱 더 일반화하였으며 기존의 방식보다 더 인간의 classification 방식에 더 가까운 방식을 제안하였음. 우리는 앞으로도 ㄹew-shot classifciation에서 우리의 방법이 baseline으로 사용되어 더 일반화된 메타 능력 측정 방식이 되기를 기대함.}

% 맨 마지막 두 문단은 어감이 이상해서 추후 수정하게Ttmqslek.
% 이 다음, GBML에 대한 조금 더 자세한 설명, 이후 특별한 feature가 task sampling procedure라는 것을 언급.
%\jh{GBML aims to train a meta-learner, \ie an initialization point which can adapt unseen classes using a few examples within a few gradient steps. \nj{For this}, GBML adopts a special procedure \nj{that}}
%\jh{\nj{samples a random} task from a task distribution and the classes in the sampled task
%are randomly assigned to labels (e.g., an integer). %This \nj{endows} a label \nj{an interesting characteristic}. Each label does not represent a specific class. This leads to a equivalence \nj{?? . That is, no label is more important, or different
%than others. 
%num-way에 대한 설명, 근데 이게 좀 사람들이 듣기 어렵지 않나. 해당 way보다 더 많은 way로 학습할 수 있는

% \jh{First feature of this equivalence implies that we do not need same number of neurons to the number of labels of episode. This is because as every label is no more different than other}
% \jh{따라서 우리는 다음과 같은 구조를 생각할 수 있다. 각 episode를 풀 때에, 여러개의 output index 중에서 해당 task를 풀 output index를 랜덤하게 정하게 되면 어차피 label간의 equivalence가 보장되기 때문에, output dimension이 task label set보다 더 커도 상관이 없다는 것이다. 따라서 우리는 해당 사실을 바탕으로 기존의 episode-dependent meta learning task를 조금 더 일반화 할 수 있다. Fig.\ref{fig:our_architecture}처럼. 즉 classifier가 존재하는 경우, classifier set의 task label set으로 전사하는 과정을 넣음으로서 meta-learning architecture에 조금 더 유연성을 추가할 수 있다.}
% \jh{task-label set보다 많은 classifier neuron의 존재는 단순히 모델의 유연성 뿐 아니라, (해당 way보다 더 많은 way로 학습할 수 있다 라는 말을 하고싶음). 기존 task보다 더 큰 set, 다시 말해, classify해야 하는 class의 수가 더 많은 set도 풀 수 있다는 것을 의미한다. 기존의 episode-dependent는 stream of episodes로 학습하기 때문에, 각 task의 label set의 크기가 같았고, 따라서 해당 model은 episode와 같은 shape를 가진 task만 풀 수 있었다. 하지만 이러한 label equivalence에 대한 분석을 바탕으로 우리는 최초로 fixed-way (\ie task label의 크기가 같은) 문제만 풀 수 있는 episode-dependent 모델이 아닌 any-way 문제를 풀 수 있는 episode-dependent 모델을 제안한다.}
% % Because the real class has no fixed relation to the assigned label in a task, every label is equivalent to others in that the assignment of a class to different labels should not make any difference. 
% This feature is unique to GBML. Although GBML relies on the same framework as supervised learning, only GBML has this feature which makes it special. In this paper, we focus on this equivalence among labels and seek how to use this specific feature and argue that this feature can be used as a double-edged sword boosting the performance of GBML if properly utilized.}
\begin{figure*}[h!]
\centering
\includegraphics[width=\textwidth]{figure/realfinalalgorithmfinal-1-2.pdf}
\caption{\jh{Overall structure of our proposed algorithm. The upper section delineates the implementation of the any-way meta-learning, while the lower section highlights the integration of semantic class information. It's noteworthy that the encoder $f$ is not updated with $g_s$ during the inner loop.}}

\label{fig:our_architecture}
\end{figure*}

% \jh{The equivalence \nj{among labels} itself can give some flexibility to the model. For example, it could fit any type of class. That is, once it is trained well, it could fit any dataset. \nj{*** fit dataset ***의 의미가 모호함.} \cite{asd} shows that it not only fits the dataset but also leads to good performance. In this paper, we focus on this flexibility, we aim to give more flexibility to GBML by \nj{exploiting} this equivalence. \nj{First, we} propose a \nj{GBML method that} can classify any number of classes. That is, suppose the model is trained to classify 5 classes, traditional GBML, or any supervised learning and its variants, can only classify 5 classes. However, our method can classify any number of classes more than 5 without sacrificing performance. Also, we would introduce some masking techniques that can improve the performance and dynamics of the model based on our algorithm.}

% 둘째로, 이러한 label equivalence한 성질은 문자 그대로 label이 한 class를 대표하지 않는다는 것을 알 수 있는데, 이는 'learning from data'라는 가장 기본적인 컨셉에서, data의 정보량이 episode-dependent한 학습 방법론은 class-label matching이 가능한 supervised learning보다 떨어진다는 것을 의미한다. 개는 첫번째 인덱스에 올 수 도 있고, 두번째 인덱스에 올 수 도 있기 때문이다.

% 본 페이퍼에서는, 메타 러닝에서 이러한 현상이 성능에 장애가 될 수 있음을 처음으로 제안하고, Episode-dependent한 학습 방법론에 label information을 추가적으로 넣음으로서 이를 해결하고 성능을 전반적으로 올릴 수 있는 단순하지만, 추가적인 연산량이 많이 필요하지 않은 방법론을 처음으로 제안한다. 또한 이러한 supervised learning 성질을 추가함으로서, 우리는 episode-dependent meta learning에 supervised learning에서 자주 사용하는 trick을 쓸 수 있음을 보인다. 이는 supervised learning에서 자주 사용하는 기법 중 하나인 mixup을 사용하여 성능을 전반적으로 올림으로서 증명한다. 또한, label information을 사용한다고 하더라도 성능 개선이 dataset-specific하게 이루어지는 것이 아님을 본 페이퍼에서 증명한다.

% 반면, episode independent한 method의 경우 위에서 언급했던 문제와는 관련이 없다고 이야기 할 수 있는데, 왜냐하면 해당 현상은 전부 episode를 통해서 학습함으로서 벌어지기 때문이다. 따라서 episode를 통해서 sampling하지 않는 방법들 \cite{rethinking..}은 이와 관련이 없다. 하지만, 반면에 이러한 방법론들은 문제를 해결하는 것을 단순히 양적으로 encoder를 좋게 만든다던가, few-shot example의 양을 늘림으로서 해결하는 등의 meta-learning이라는 essense와는 떨어지는 것이 사실이다. 따라서, 본 논문에서는 아주 작은 수준의 추가적인 조건만 가지고 episodic 한 성질을 inject하고, 성능을 드라마틱하게 개선할 수 있는 방법을 제안한다.
% % \jh{\nj{secondly,} we will address that this equivalence naturally leads to a lack of information. That is, as all labels are equivalent, the assigned 
% target label lacks information about the exact class of the sample in question. For example, in any traditional supervised learning, each class represents
% a specific class and a model can learn about the class by classifying each input to the corresponding class. However, in GBML, as \nj{all the labels are} equivalent, the model cannot learn about the specific class to which a sample belongs. %In other words, the model lacks information on that class compared to supervised learning. 
% In this paper, we are to solve this problem by injecting label information into the model. Our algorithm does not require specific architecture or a huge computational \nj{overhead}. Also, injecting label information means that we can make use of the techniques utilized in supervised learning such as \cite{mixup, } in GBML. In this paper, we propose to apply Mixup \cite{mixup} technique to GBML and show that it improves the performance. It suggests that some techniques of supervised learning can be applied to GBML within our framework.}
% % 본 페이퍼에서, 우리는 가장 대표적인 episode-dependent paper인 MAML에 우리의 방법론을 적용해보고, 우리의 방법론이 잘 work한다는 것을 보일 것이다. 그리고 우리의 방법론이 몇가지의 가벼운 조건만 있으면 우리의 방법론이 episode-independent에도 적용 가능하다는 것을 선보일 것이다.

% 우리 페이퍼가 전반적으로 넓은 분야를 지칭하기 때문에, 해당 분야에 대한 설명을 Related Work에서 하도록 한다.
\section{Related Works}
% 메타 러닝에 대한 간략한 설명, 무엇을 이야기해야 하냐면, 이것의 목표는 뭐고, 대부분은 무엇무엇을 다룬다. 어떤 모델이 메타 러닝 능력을 가지고 있다고 이야기하기 위해서라면, 일반적으로 그 모델은 몇 개의 샘플들만 가지고 그 샘플들간의 관계에 대해서 알 수 있어야 할 것이다. 이런 개념을 \cite{labelequiv:metadefinition}에서 정의하였으며, 다음과 같이 meta-learning ability를 지속적으로 체크하며 직접적으로 성능을 발전시키는 매커니즘을 \cite{label_equiv:matchingnetwork}에서 제안하였다. ㅐ해당 논ㄴ문에서는 각각 작업을 풀기 위해서 작은 데이터셋 (support set)을 사용하는 episodic training 방식을 도입하였다. 

% \subsection{Episode-based Meta Learning}
% Matching network \cite{label_equiv:matchingnetwork}을 필두로 하여, 현재의 대부분의 메타 러닝 모델은 episode-based로 학습한다. \cite{nonepisodic:transferminitocars}가 정리한 것처럼, 대부분의 episode-based model의 경우 meta learning 능력을 얻기 위한 meta-training stage, 그리고 이 성능을 테스트하기 위한 meta-test stage로 이루어져 있으며, 크게 두 가지의 방식의 접근법이 사용된다

% 첫째 방법론은 Gradient-Based Meta Learning이다(GBML) Model-Agnostic Meta-learning model(MAML) \cite{episodic:GBML:MAML}이 첫 번째로 제안하였으며 해당 방식은 support set을 inner loop에서 gradient descent로 학습하고, 얼마나 잘 optimized 되어 있는지를 query set에서 학습한 후, 더 support set을 잘 학습하도록 outer loop에서 meta-initialization paremeter를 학습시킨다. GBML이라고 불리는 연구는 많이 이루어지고 있으며 \cite{episodic:GBML:ANIL, episodic:GBML:imaml, episodic:GBML:warpgrad, episodic:GBML:Reptile} 이때 inner loop에서의 성능을 계산할 때 Hessian을 계산해야 하기 때문에 이것의 계산량을 latent vector를 이용하여 줄이거나 \cite{episodic:GBML:LEO}, 혹은 Hessian을 이용하지 않는 방법들이 연구되어왔다 \cite{episodic:GBML:ANIL, episodic:GBML:Reptile, episodic:GBML:MAML}. 
% 두번째 방법론은 Metric-Based MEta Learning (MBML)으로서, 좋은 encoder를 학습하여 feature 단에서 잘 clustering 되도록 하는 것이 목표이다.  \cite{label_equiv:matchingnetwork} 등에서 시작하였으며, GBML 이전부터 사용되던 방법론이다. Protonet \cite{episodic:MBML:Protonet}이 기초가 되는 아키텍쳐를 제안하였으며, 이는 support set에서 encoder를 뽑은 다음에 clustering한 다음, query set에서 해당하는 class와의 거리를 줄이는 것을 방법론으로 한다. 이때의 clustering은 Euclidean space에서 했는데, 이것은 Hyperbolic space와 같은 다른 Riemannian 공간에서 해도 잘 된다. \cite{episodic:MBML:Hyperbolic}
% \subsection{Transfer-Learning Based few shot learning}
% MBML은 때로는 단순히 좋은 encoder를 unsupervised 단에서 잘 만들도록 하는 문제에서 지표를 측정하는 metric으로 분류되기도 한다. \cite{nonepisodic:transferisbetter}가 제안하였듯이, 이 문제를 meta-learnin 문제가 아닌 few-shot learning으로 보게 된다면 encoder를 잘 학습하고, 그리고 transfer learning에서 dataset scarcity problem으로서 볼 수 있게 된다. 이러한 방법들의 공통점은 결국 unsupervised manner에서 문제를 풀어야 한다는 것인데, 이 문제를 해결하기 위하여 이들은 label에 의존하지 않고 새로운 문제를 만들어 그 문제를 풀도록 한다. 이들은 보통 두 개의 같지만 다르게 표현된 이미지를 encoder가 같게 만듦으로서 문제를 푼다 \cite{label_equiv:ssl:byol,label_equiv:ssl:simclr}. 예를 들어서 self-supervised learning의 경우 이미지를 마구 뒤섞어 직소 퍼즐을 맞추는 식으로 바꾸거나 \cite{label_eqiv:ssl:jigsaw}, 혹은 사진을 뒤집고, 회전하여도 같은 이미지라고 보게 하거나, \cite{label_eqiv:ssl:rot} 결국에 \cite{label_eqiv:ssl:rot,label_equiv:ssl:byol}. 이렇게 학습시킨 후, 실제 데이터셋으로 들어가기 전 모델은 classifier를 학습하는 transfer learning을 거친다. 

% 여기서 주목할만한 점은, 이 문제를 transfer learning으로 바꾸어서 풀게 되더라도 결국에는 label information이 없기 때문에 어떤 task를 풀기 위해서 절대적인 class를 represent하는 semantic class supervised learning을 거친다는 점이다. 이는 task로만 학습해서 develop하는 episodic learning과는 대비된다. episodic learning에서는 semantic class가 사용되지 않고, 각 task에서 numberic label만 랜덤하게 배정된다.
% \subsection{Meta Learning} \yr: 이렇게 하는 건 별로야? ㅎㅁㅎ
\textbf{Meta-learning}, briefly described, is about the study of algorithms that learn from data. As defined in \cite{labelequiv:metadefinition}, the fundamental concept of meta-learning is the model's capability to discern relationships among a limited number of samples.
%The crux of meta-learning is the capability of a model to discern relationships between a limited number of samples. This fundamental concept of meta-learning was defined in \cite{labelequiv:metadefinition}. 
Furthermore, mechanisms were proposed in \cite{label_equiv:matchingnetwork} to continuously verify meta-learning abilities while simultaneously enhancing performance.
%mechanisms that continuously verify meta-learning abilities while directly enhancing performance were proposed in \cite{label_equiv:matchingnetwork}. The latter introduced the episodic training method, employing a small dataset, termed the support set, to tackle each task.

Inspired by the Matching network \cite{label_equiv:matchingnetwork}, a majority of current meta-learning models adopt an episode-based training approach. As summarized by \citet{nonepisodic:transferminitocars}, these episode-based models usually consist of two stages: the meta-training stage to acquire meta-learning capabilities, and the meta-test stage to evaluate their performance. There are two lines of research in this area:

\textbf{Gradient-Based Meta Learning (GBML)} %This methodology,
first introduced by the Model-Agnostic Meta-learning model (MAML) \cite{episodic:GBML:MAML}, involves training the support set in the inner loop through gradient descent. The degree of optimization is then evaluated on the query set, followed by the outer loop refining the meta-initialization parameter to enhance learning on the support set. %The degree of optimization is then gauged on the query set, after which the outer loop refines the meta-initialization parameter to improve learning on the support set. 
There's a growing body of research in GBML
%, as evidenced by works 
such as \cite{episodic:GBML:ANIL, episodic:GBML:imaml, episodic:GBML:warpgrad, episodic:GBML:Reptile}. Notably, due to the computational demands of \jh{caluculating the Hessain} \nj{in the inner loop}, %performance computation, 
methods that utilize latent vectors to reduce this computation %to reduce this computation using latent vectors 
have been explored \cite{episodic:GBML:LEO}. Other studies have put forward the idea of % Other research efforts have proposed 
avoiding the use of the Hessian altogether \cite{episodic:GBML:ANIL, episodic:GBML:Reptile, episodic:GBML:MAML}.

\textbf{Metric-Based Meta Learning (MBML)} is focused on training an effective encoder that can perform clustering at the feature level. This approach, which predates GBML, was established by initiatives such as \cite{label_equiv:matchingnetwork}. 
%Initiatives such as \cite{label_equiv:matchingnetwork} set the foundation for this methodology, which predates GBML. 
A foundational architecture for this was proposed by Protonet \cite{episodic:MBML:Protonet}. The method involves extracting features for the inputs in the support set by an encoder, clustering them, and then minimizing the distance between the query set and the corresponding class. Although clustering is conventionally conducted in the Euclidean space, it has been shown to be more effective in other Riemannian spaces like the Hyperbolic space \cite{episodic:MBML:Hyperbolic}.
% \jh{textbf{Domain Generalization} Domain generalization is type of research which seeks a robust model with different environmenment \cite{DG:Gunbon0,DG:Gunbon1}. Our any-way meta learning also deals with same situation. As shown in Table.~\ref{table:intuition}, episodes with different cardinality can be treated as different distribution. Since our goal is to implement a model with true-meta-learning-ability(a model which has meta-learning ability within any task cardinality), we also deal with domain generalization. The conventional wisdom is that although we can increase overall performance of given task, the performance of each dataset decreased compared to the episode-sole-trained dataset \cite{DG:wisdom0,DG:wisdom1,DG:wisdom2}. In this paper, we argue that by using label equivalence, we show this wisdom might not be hold in meta learning area.. Which means, It not only adapts well any-way, but also actually perform better than fixed-way training.}

\textbf{Domain Generalization} field focuses on designing
% \jh{\textbf{Domain Generalization:} The field of domain generalization aims to design 
models that exhibit robustness across diverse environments~\cite{DG:Gunbon0,DG:Gunbon1}. Our approach to ``any-way" meta-learning aligns with this objective. As illustrated in Table~\ref{table:intuition}, episodes of varying cardinality can be perceived as \nj{from} distinct distributions. In pursuing a model endowed with a genuine meta-learning capability— capable of adapting to any task cardinality— we inevitably confront the challenges posed by domain generalization.

The \jh{conventional} understanding suggests that while it is possible to enhance the overall performance on a set of tasks, the performance on individual datasets might decline %may suffer
when compared to models trained exclusively %solely 
on a specific episode~\cite{DG:wisdom0,DG:wisdom1,DG:wisdom2}. In this paper, we challenge this established notion. %conventional wisdom.
Through the integration of label equivalence, we posit that this trade-off might not be as significant %pronounced 
within the meta-learning realm. Our results \nj{exhibit} not only the seamless adaptability to ``any-way" tasks but also superior performance compared to fixed-way training paradigms.

% \jh{\textbf{Transfer-Learning Based Few-shot Learning}
% Sometimes, MBML is categorized as a metric to \yr{evaluate} the quality of an encoder crafted in an unsupervised manner. As suggested by \cite{nonepisodic:transferisbetter}, if this challenge is viewed not as a meta-learning problem but as a few-shot learning task, the focus shifts to proficient encoder training. Furthermore, it sees the challenge from a transfer learning perspective, addressing the dataset scarcity problem. A common theme across these methods is solving the problem in an unsupervised fashion. They construct novel problems that don't rely on labels. For instance, in self-supervised learning, strategies like rearranging images to match jigsaw puzzles \cite{label_eqiv:ssl:jigsaw}, or ensuring encoders perceive flipped or rotated images as identical \cite{label_eqiv:ssl:rot, label_equiv:ssl:byol}, are commonly employed. After such training, the model undergoes transfer learning by training a classifier before interacting with the actual dataset.
% }
% \jh{An essential observation here is that even when reframing the problem through transfer learning, there's an inherent dependency on semantic class supervised learning due to the absence of label information. 하지만, few-shot classification을 할 때에 있어서는 그러한 과정이 생략된다. This is in contrast to episodic learning where semantic classes aren't used, and numeric labels are randomly assigned for each task.}
% \jh{또 하나의 중요한 점은 few-shot classification 능력을 측정할 때에는 동일한 sampling mechanism을 가진다는 점이다.}
% \section{Preliminaries}
% \begin{itemize}
%     \item Gradient-based Meta Learning
%     \item inner loop 설명
%     \item fixed-way 설명
%     \item equivalence 다시 설명 
%     \item true label, meta label 정의 필요, task-label, class-labelk
% \end{itemize}
\section{Method}
\subsection{Conventional (Fixed-way) Meta-Learning}
% \begin{figure*}[h!]
% \centering
% \includegraphics[width=\textwidth]{figure/sampling_algorithm.png}
% \caption{\jh{sampling algorithm. By doing so, every classifier can backpropacate in every episode(task)}}
% \label{fig:masking? procedure}
% \end{figure*}

To evaluate the few-shot learning ability within a given distribution, we define a process of task sampling. Initially, 
%Firstly, 
we select a label set from the class pool of a larger mother dataset  $\mathcal{D}$. In conventional meta-learning, the cardinality $N$ of the selected label set remains
%is
fixed during the entire training procedure. We then sample data from this set, each matching the size of the support set (number of shots).
Each task $\tau$ consists of a set $\mathcal{X}$ of tuples $\{(x,y)\}$ sampled from $\mathcal{D}$, where $x$ denotes the input data and $y$ represents the numeric label rather than the semantic class from the dataset (Fig. \ref{fig:task_sampling}).  Moreover, $\mathcal{X}$ is divided into a support set $\mathcal{X}_s$ and a query set $\mathcal{X}_q$ for task learning and performance evaluation, respectively. % 이렇게 쓰는 건 어떨까?
% 기존: Moreover, $\mathcal{X}$ is partitioned into a support set $\mathcal{X}_s$ used for learning the task and a query set $\mathcal{X}_q$ \yr{for performance evaluation.} <- to evaluate its performance. ( for ~, 구조 맞추고 싶어서 ㅎㅁㅎ)
To do so, a model with encoder $f$ and $N$-way classifier $g_{f}$\footnote{Subscript $f$ denotes `fixed-way' while $a$ is used for `any-way'. } learns from the support set, then \jh{evaluates} its performance with the query set.

% \jh{우선, 우리의 모델은 Output dimension $O$의 크기가 task label set의 크기보다 크기 때문에 매번 Task를 뽑았을 때 $O$중 Task label에 해당하는 $|T|$개의 label set을 selection하는 과정을 가져야 한다. 또한 equivalence를 유지하기 위해서 매번 새로운 Episode마다 $O$중 $|T|$를 sampling하는 sampling combination $s$는 $C(|O|,|T|)$에서 random하게 뽑는다.}
\begin{algorithm*}
\caption{Any-Way Meta-Learning}\label{alg:AnyWayMetaLearning}
\begin{algorithmic}

\Require Output dimension $O$ %, number of class in a task $N$} %classifier nodes set $C$, task set $t$
\Ensure Trained model $\theta$
\State Initialize model parameters $\theta$
\For{episode $= 1$ to $K$}
    \State Sample a random number $N \le O$
    \State Sample a task $\tau$ with $N$ number of classes and set $J = \lfloor O/N \rfloor$.
    \State Generate $S = \{s_j\}, j \in [J]$ where $s_j \in \mathbb{R}^N$ is a random assignment, non-overlapping with $s_i, \ \forall i\neq j$ 
    %\State Calculate $S(g(f(x_{\text{support}})))$, which has dimension $(\lfloor\frac{O}{|T|}\rfloor, |T|)$
    \State Calculate any-way loss $L_a = \sum_{j=1}^J L(S_j(g_a(f(x_{\text{support}}))),y)$ and update the model parameters $\theta$ using $L_a$
    \State Calculate logit for the query set: logit $=\sum_{j=1}^J S_j(g_a(f(x_{\text{query}})))$
    %\State Initialize logit $\in \mathbb{R}^N$ to zero for performance testing with query set
    %\State $\text{logit} \in \mathbb{R}^N$
    %\For{each $j$ in $J$}
    %    \State $\text{logit} \mathrel{+}= S_j(g(f(x_{\text{query}})))$
    %\EndFor
    \State Calculate accuracy, and update meta-parameters based on logit and the actual labels of the query set
\EndFor
\State \textbf{return} trained model with parameters $\theta$

\end{algorithmic}
\end{algorithm*}

\subsection{Any-way Meta-Learning}\label{sec:anywaymetalearning}
As mentioned earlier, there exists an equivalence among numeric labels. % between은 둘 사이 느낌이고, among은 좀 더 여럿. numeric labels면 among이 맞지 않을까?
We can leverage
% exploit은 뭔가 착취?하는 느낌 
these characteristics to enhance the meta-ability of the model.

% a task-classifier $g_t$ that classifies these features into task-labels, and a sampling procedure, denoted as $s$. The function of $s$ is to identify the index in the output $O$ of $g_t$ that corresponds to the task-label. Mathematically, this is expressed as $s(g_t(f(x)))= { g(f(x))[i] : i \in s}$, where $s$ belongs to the set of combinations $C(|O|,|T|)$. Here, $C(|O|,|T|)$ refers to all possible ways of selecting $|T|$ elements from a set of $|O|$ elements. Usually, the output dimension of $g_t$ matches the task-label space set, implying that $s$ is a complete set selection. 
% \jh{In this paper, we introduce two significant modifications to the conventional approach. Firstly, we incorporate task cardinality as a variable in our task sampling process during training. Secondly, we expand our model by integrating a class-classifier, $g_c$, specifically designed for class-label classification.}

Firstly, our model's output dimension, $O$, is set to be greater than the maximum cardinality of the tasks' label sets, implying that each time we sample a task $\tau_i$ with $N_i$ classes, we must assign $N_i$ labels among $O$ output nodes. % corresponding to the task label. This assignment should be random to maintain equivalence.

%for each new episode, we randomly select labels. Which means, we draw a random sampling combination $s$ from $C(|O|,|T|)$, where $s$ indicates the $|T|$ labels from $O$.}
With our proposed configuration % proposed configuration? 어떤 걸 의미하는거지?
, we maintain output-to-class assignments during each episode. This not only determines a unique and unchanging pathway for selecting numeric labels but also restricts gradient updates to occur exclusively along this fixed way, thereby accommodating the execution of standard meta-learning tasks. Unlike traditional methods that fix the cardinality of the task when conducting task sampling, our approach also involves sampling the cardinality. % yr: 문장 호흡이 너무 길어서 나눔
This enables our model to handle tasks with any cardinality, hence fostering a cardinality-independent learning environment, \ie any-way meta-learning. 
%\yr{learning environment that is independent of cardinality, \ie any-way meta-learning.}
However, a potential problem arises when the output dimensionality $O$ significantly exceeds the expected task cardinality. In this case, the chance of each classifier node being selected %to classify the task 
in each episode approaches zero. Consequently, classifier nodes might undergo less training than the encoder throughout the training procedure, resulting in an underfit.
%at $\frac{|O|}{\mathbf{E}(|T|)}$. Should this probability be low, classifier nodes might undergo less training than the encoder throughout the learning process. In this paper, we introduce a simple yet effective remedy to counteract this problem.
To remedy this, we assign a set of output nodes to a numeric label. We assign a numeric label to (almost) all nodes for classification, rather than assigning a numeric label to just one output node and ignoring the unassigned nodes.
% \yr{rather than? 을 말한 이유 잘 모르겠음. to clarify? 아니면 기존엔 저렇게 해?} Rather than assigning \nj{a numeric label to only one output node} %three output nodes \nj{to the three numeric labels} 
%and ignoring the \nj{unassigned \yr{nodes}}, we propose assigning a numeric label to all nodes for classification.

Let $f(x)$ be a feature map for the input $x$ and $g_a(\cdot)$ be an any-way classifier consisting of $O$ output nodes. For a task with $N$ classes, we can constitute $J = \lfloor O/N \rfloor$ assignments $s_j \in \mathbb{R}^N$, $j \in [J]$ whose $i$-th element corresponds to the output node for the $i$-th numeric class label\footnote{For example, consider a scenario with eight classifier output nodes and three classes in a task, thus, three numeric labels. In this case, it can be $s_1 = [3,5,2]^T$, $s_2 = [7,4,8]^T$ and nodes 1 and 6 are unassigned. }.  Here, $\lfloor \cdot \rfloor$ is the floor operation, and $[n]$ denotes a set of natural numbers from $1$ to $n$. We then aggregate the loss for each assignment to constitute our any-way loss $L_a$ as
\begin{equation}
L_{a} = \sum_{j=1}^J L(S_j(g_a(f(x))),y).
\label{eq:any-way-loss}
\end{equation}
Here, $S_j: \mathbb{R}^O \rightarrow \mathbb{R}^N$ uses the vector $s_j$ as an index function to extract $N$ elements, \eg $S_1(v)$ with $s_1 = [3,5,2]^T$ outputs $[v_3, v_5, v_2]^T$.    As different assignments do not overlap, \ie $\text{set}(s_i) \ \cap \ \text{set}(s_j) = \phi$ where $\text{set}(v)$ is the set consisting of elements of the  vector $v$, this additional assignment procedure %merely selects different output nodes, it 
requires no additional computation. \ie no additional backpropagation nor inference.  

%We can denote the first assignment set as $s_1$ and the second as $s_2$\footnote{Here, $s_1$ and $s_2$ are ordered selections of $N$ output nodes whose $i$-th element being assigned to numeric label $i$, e.g. $s_1 = \{3,5,2\}$ and $s_2 = \{4,1,6\}$.} in this senario. Consequently, the original loss becomes $L(s_2(g(f(x))),y) + L(s_1(g(f(x))),y)$. As this additional procedure merely selects different output nodes, it requires no additional computation. \ie no additional backpropagation nor inference.

% \jh{To formalize this mathematically, we construct a combination bundle set $S$. Let $S = \{s_0, \dots, s_i\}$, where $i = \lfloor{\frac{|O|}{|T|}}\rfloor$. Also, the elements of $S$ are disjoint, i.e., $s_i \cap s_j = \emptyset \quad \forall i,j$. At every iteration, by drawing a different combination set $S$ and constructing $S \circ g \circ f$, all output nodes can have task labels assigned, thereby solving the corresponding problem.}

% \jh{As a result, upon completion of training, the model can possess multiple independent and equivalent node sets capable of classifying a given task, all of which can be produced simultaneously in a single round of training. As each set of node can classify incoming query labels, we regard these node sets are independent models. By ensembling these nodes, we can not only naturally extend to any-way scenarios but also induce performance improvement. While it is generally expected that any-way learning should have a less restrictive upper bound than fixed-way learning, the ensemble effect previously mentioned can surfass the fixed way model's performance.}

%\jh{이렇게 되면 매 task마다 task label을 뽑아내는 forward path가 고정되고, 여기에서 우리는 일반적인 Meta learning Task를 수행하면 된다. 이때 $|O| > |T|$라는 조건만 만족한다면 앞서의 과정을 수행할 수 있기 때문에, 우리는 어떤 number of task label set이든 다 학습할 수 있다. 따라서 우리의 방법론은 Task sampling을 할 때에 way를 fix하는 기존의 방법과는 다르게 $|O|<|T|$를 만족하는 way도 마찬가지로 랜덤하게 샘플링하여 어떤 way등 agnostic하게 학습할 수 있도록 하자. (고쳐야 할 부분 - way라는 단어의 대용어? number of task label은 너무 김)}

% \jh{해당 아이디어는 어떤 task-label size agnositc하게 동작하게 하지만, 반면에 $|O| >> {EX}(|T|)$의 경우에는 문제가 생길 수 있음을 알 수 있다. 왜냐하면 매 episode(task)마다 각각의 classifier neuron이 task를 classify하도록 선택될 확률은 $\frac{|O|}{EX(|T|)}$이며, 만약 이 확률이 낮다면 학습이 진행되는 동안 classifier neuron 들은 training이 encoder에 비해서 상대적으로 덜 train될 것이기 때문이다. 따라서 본 페이퍼에서는 이를 해결할 수 있는 간단하지만 효과적인 방법을 제시한다.}

% \jh{이를 하기 위해서 우리는 set of task class label을 할당한다. 예를 들어 classifier neuron의 개수가 6개이고, task set이 3개라고 가정하자. 그러면 classifer neuron중 3개만 선택하고 나머지는 무시하는 것이 아닌, 이중 나머지 3개의 neuron에도 똑같이 task-label을 assign해서 문제를 classify하자는 것이다. 이중 전자의 selection을 $s_1$, 후자의 selection set을 $s_2$라고 하면, 기존의 loss는 $L(s_2(g(f(x))),y) + L(s_1(g(f(x))),y)$로 바뀌며, 이는 똑같은 output 값들을 서로 다르게 select한 것일 뿐이니 추가적인 계산이 전혀 필요하지 않다. 이를 수학적으로 엄밀하게 표현한다면. \jh{combination} bundle set S를 구성한다. $S = \{s_0, \dots, s_i\}$ 고, $i = \frac{|O|}{|T|}$ 이다. 또한, $S$의 element들은 서로 disjoint이다. \ie $s_i \cap s_j = \emptyset \forall i,j$ 이다. 매 iteration마다 combination set S를 서로 다르게 뽑아 $S\circ g \circ f$로 만들면 모든 output neuron에 대해서 task label이 assign되고, 따라서 해당하는 문제를 해결할 수 있다. }

% \jh{결과적으로 학습이 완료되었을 때에, 모델은 주어진 task를 classify할 수 있는 \textbf{여러개의 독립적이고 equivalent한 neuron set}을 가질 수 있게 되며, 이들은 한번의 training으로 동시에 다 만들어 낼 수 있다. 우리는 이런 neuron들을 ensemble함으로서 단순한 any-way로의 확장 뿐만 아니라 자연스러운 성능 향상을 이끌어 낼 수 있다. 일반적으로 any-way는 fixed-way 학습보다 더 제약조건이 부족하기 때문에 fixed-way가 upper bound가 ㄷㅣ어야 하지만, 앞서 언급하였던 앙상블 효과 때문에 오히려 성능이 더 올라가 ㄹ여지가 있다. }


\subsection{Injecting Semantic Class Information to Model}
% As previously mentioned, \yr{a limitation arises as a numeric label does not represent a particular class, hence containing less information than in conventional supervised learning. 이 문장을 어떻게 하면 맛있을까..} 
% \yr{To mitigate this issue, we append an auxiliary classifier designed to distinguish the semantic classes.} 
% %In this section, we propose a method to mitigate this issue \nj{by appending} an auxiliary classifier \yr{designed} to distinguish the semantic classes. 
While it's possible to
%Although we could 
implement any-way meta-learning through label equivalence, relying solely on numeric labels leaves out important semantic information. To overcome this, we introduce an auxiliary classifier to fill this gap.
% In the face of numeric labels' inherent limitations, we introduce an auxiliary classifier, purposefully designed to discern nuanced semantic classes. 
The overall structure is illustrated at the bottom of Fig.~\ref{fig:our_architecture} ($g_s$).

During the inner training loop, the auxiliary classifier employs features produced by the encoder. These features originate from an initial set point defined by meta-parameters. Our main objective is to ensure that these features encapsulate %carry 준후 기준에 박탈당할 것 같긴 한데 encapsulate이 젤 적합한 단어 같아서 소매넣기 해봄 skeh rmwjdehsms dkfdk...
the semantic information of the class. To verify this, we use the auxiliary classifier to discern different semantic classes. To maintain %For the sake of 
consistency and simplicity, the encoder remains unchanged \yr{throughout} this process, allowing for a single pass through the network. The majority of computations are streamlined through a compact linear layer, ensuring minimal computational overhead.
% Unlike traditional models that only receive numeric labels of a task as targets, % in every episode, 
% our method \yr{leverages} semantic class\nj{es} as \yr{supplementary} %yr: additional 말고 supplementary는 어떤가유?
% target labels. We train an auxiliary classifier to classify these target labels at each episode. Within the inner loop, we train this auxiliary classifier using encoder features \yr{obtained} from the meta-parameters (initialization point) to classify the semantic class. \yr{Since the purpose here is to \yr{enable} the encoder features extracted at the meta-initialization point to distinguish semantic class\yr{es}, the encoder remains fixed, thus necessitating only a single feed-forward pass. 이 문장 이해가 잘 안 됨} The actual computation occurs within a small linear layer, introducing only a \yr{minimal} %slight 
% computational overhead. 

In our implementation, the semantic loss, $L_s$, which is a typical cross-entropy loss with $C (\gg N)$ class\nj{es}, is added to the original loss (either a fixed-way loss or an any-way loss in (\ref{eq:any-way-loss})) using a balancing weight $\lambda$, \ie $L_{total} = L_{original} + L_s$.
%\jh{We used obtained this semantic loss as a regularlizer in the outer loop. \ie $L_{\text{original}} + \lambda L_{semantic}$}

The inclusion of the auxiliary classifier provides a gateway to information that was previously inaccessible %untapped 
in conventional episode-based meta-learning. This unique structure, where supervised learning is inherent in the auxiliary classifier in every episode, paves the way for incorporating conventional supervised learning techniques into our approach.
As an exemplar, we employ the Mixup technique, one of the most commonly used data-augmentation techniques in supervised learning. Using the known semantic class\nj{es}, we apply this technique as it is used in supervised learning. %in a manner consistent with its traditional use in supervised learning. 
While the numeric labels stay unchanged, we blend the semantic class across tasks, generating new input-output pairings.
This practical implementation of the Mixup technique in our methodology emphasizes our central assertion: conventional supervised learning strategies can be seamlessly and effectively integrated within the meta-learning landscape.
% The addition of this auxiliary classifier \yr{enables} us to obtain information that was \yr{previously} inaccessible in the traditional episode-based meta-learning. Furthermore, \yr{given that} supervised learning takes place in the auxiliary classifier in every episode, we can leverage techniques used in supervised learning. \yr{
% In this work, we present an \nj{example} that utilizes Mixup \cite{??}, one of the most commonly used data-augmentation techniques in supervised learning, which involves blending inputs from different classes to generate new images and creating corresponding labels according to the mixing ratio, \nj{to further boost the meta-learning performance}.
% }
% \yr{With knowledge of the semantic class, we apply the mixup in the same manner as in conventional supervised learning. Here, we retain the numeric labels of a task unchanged, instead mixing the semantic class across different tasks to generate new input-output pairs.}

% By implementing the mixup \yr{technique} in our algorithm, we \yr{demonstrate} the feasibility of \yr{integrating
% } conventional supervised learning techniques into the \yr{domain} of meta-learning.

\renewcommand{\arraystretch}{1.2} % 1.5는 줄 간격을 1.5배로 늘립니다.
\begin{table*}[]
\caption{\jh{f denotes ‘fixed-way’ while a is used for ‘any-way’. For implementing 10-way-1-shot f-MAML, we encountered many corruption in this case. \ie the model does not trained. We discarded all these failed experiments and sampled three 'successed' models with different seeds for performance report. For a-MAML, we did not encountered any corruption.}}
\centering
\begin{center}
    \adjustbox{max width=1.0\textwidth}{
    \begin{tabular}{l|ccc|ccc|ccc|ccc}
    Dataset & \multicolumn{3}{c|}{miniImageNet} & \multicolumn{3}{c|}{TieredImageNet} & \multicolumn{3}{c|}{Cars} & \multicolumn{3}{c}{CUB} \\ \hline
    Num-Way                   & 3 & 5 & 10 & 3 & 5 & 10 & 3 & 5 & 10 & 3 & 5 & 10 \\ \hline
    f-MAML (1)                  & 59.56(1.48) & 46.86(0.95) & 31.91(0.50) & 46.86(0.95) & 46.76(0.58)& 46.86(0.95) & 64.02(0.95) & 49.19(0.56) & 32.84(1.91) & 46.86(0.95) & 57.35 (0.58) & 46.86(0.95) \\ 
    a-MAML (1)           &63.15 (0.52)   & 48.33(0.26)  & 31.34(0.28)   &61.92 (0.26) & 47.43(0.28)  & 31.62(0.18)   &64.72 (0.24)   & 50.24(0.23)  & 33.79(0.26)   &72.21 (0.19)   &58.85(0.33)   &41.66(0.42)    \\
    + semantic (1) &64.25 (1.42)   & 49.16(0.83)  & 32.28(0.87)   &62.80 (0.92)   & 48.33(1.02)  & 32.41(0.80)   & 65.08(0.94)  & 50.93(0.89)  & 34.04(0.33)   & 72.63(0.22)  & 59.27(0.45)  & 41.82(0.96)   \\ \hline 
    f-MAML (5)                  & 76.64(0.41)  & 65.41 (1.04)  & 51.18(0.27)   & -  & 67.69 (0.23)  & -   & 79.54(0.85)  & 66.59 (1.22)  & 50.05(1.38)   & -  & 74.71 (0.40)  & -   \\
    a-MAML (5)          & 79.06(0.38)  & 66.73(0.42)  &  50.28(0.49)  & 79.30(0.47)  & 68.10(0.63)  & 52.19(0.78)   & 81.55(0.49)  & 70.42(0.88)  &  52.67(2.23)  & 84.77  &  75.62 & 60.8   \\
    + semantic (5) & 78.89(0.42)  & 66.91(0.10)  & 50.16(0.0.06)   & 79.08(1.26)  & 67.40(1.29)  & 51.12(1.73)   & 81.65(0.19)  & 70.56(0.12)  & 52.125(0.78)   & 84.77 (0.15) & 75.36(0.12)  & 59.76 (0.26) 
    \end{tabular}
    }
\end{center}
\label{table:GBML_MAIN}
\end{table*}



% 앞서 언급했던 것처럼 task sampling을 할 때 task-label이 특정 class를 represent하지 않기 때문에, 정보량이라는 관점에서는 다른 supervised learning보다 적은 정보를 retrieve한다느 ㄴ한계가 있다. 본 섹션에서는 이러한 문제를 해결하는 방법을 제시한다. 이를 해결하기 위해서 우리는 auxiliary classifier를 추가로 달아 해당 classifier로 하여금 class의 label을 구분하도록 한다. Fig.~\ref{fig:our_architecture}에 전체적인 구조가 설명되어 있다. 기존에 매 episode마다 모델이 task label만을 target으로 제공받는데 반해, 우리의 method는 class label또한 target label로 사용한다. 이때 target label을 classify할 수 있는 auxiliary classifier를 추가로 달아서 학습시킨다. 매 episode마다. inner loop에서 meta-parameter(initialization point)에서 뽑아낸 encoder feature를 사용해 class-label을 classify하도록 auxiliary classifier를 학습한다. 이때의 목적은 meta-initialization point에서 뽑아낸 encoder feature가 class-label을 구분할 수 있도록 하는 것이기 때문에 encoder는 고정되고, 따라서 결국에 feedforward를 한번만 하면 된다. 실제로 계산이 이루어지는 것은 작은 linear layer이다. 따라서 계산량이 크지 않다.

% 이러한 auxiliary classifier를 추가함으로서, 우리는 기존의 episode-based meta learning에서는 얻을 수 없었던 정보를 얻을 수 있다. 또한 매 episode마다 auxiliary classifier에서는 supervised learning이 이루어진다고 할 수 있기 때문에 이러한 성질을 활용하여 우리는 supervised learning에서 사용되는 technique를 사용할 수 있다. 본 work에서는 가장 많이 쓰이는 technique중 하나인 mixup을 활용한 알고리즘을 제시한다.

% \begin{figure*}[h!]
% \centering
% \includegraphics[width=\textwidth]{figure/mixup.jpeg}\label{fig:mixup}
% \caption{\jh{Mixup procedure, fix the task label, then mix different class-labeled images. So mixup image maintains same task idx}}
% % 같은 class 다른 색 강아지 rotation, flip?
% \label{fig:tasksampling}
% \end{figure*}

% Mixup은 다음과 같이 이루어진다고 할 수 있다. 첫째, 서로 다른 task를 sampling하였다고 하더라고 그 두 task는 같은 개수의 label set을 가지고 있고, 또 서로 다른 true label을 가지고 있다. (해당 true label, label은 따로 notation 정리가 필요함) 이때 같은 label을 가지고 있는 서로 다른 task의 image들을 mixup한다. 따라서 이들은 여전히 기존의 inner loop 관점에서는 같은 optimization task를 풀고 있지만, true-label을 풀어야 하는 auxiliary classifier 입장에서는 mixup 문제가 된다. \ref{fig:mixup} 

\section{Experiment}
\subsection{Implementation Details}
\textbf{Datasets} We \yr{evaluated} our methodology on a diverse \yr{range} of datasets. The general datasets, denoted as `G', such as miniImageNet \cite{dataset:miniImageNet} and \yr{T}ieredImageNet \cite{dataset:tieredimagenet} (subsets of the more extensive ImageNet \cite{dataset:imagenet} with 100 and 400 classes respectively) serve as versatile bases for broader tasks. In contrast, the Cars \cite{dataset:cars} and CUB \cite{dataset:cub} datasets, representing specific datasets or `S', are widely \yr{used for} fine-grained image classification evaluations \yr{due to} their focus on closely related objects with \yr{subtle} variations. 
By utilizing these datasets, we are able to comprehensively evaluate the performance of our methodology across different task spectrums. % The adoption of these datasets facilitates a nuanced assessment of our methodology across different task spectrums.

\textbf{Task sampling} For each episode's initiation, the task cardinality was sampled from \nj{\{3, 5, 7, 9\}}. \yr{While \nj{tests} were} %Despite tests being
conducted on 3-way, 5-way, and 10-way, the 10-way cardinality was excluded from our sampling pool.

\textbf{Hyperparameters} In line with \cite{episodic:GBML:BOIL}, our experiments involved sampling 60,000 episodes. We adopted the 4-conv architecture as detailed in \cite{dataset:miniImageNet}. The learning rates were set at 0.5 for the inner loop and 0.001 for the outer loop. Depending on the shot type, the $\lambda$ values were adjusted: 0.1 for 1-shot and 0.5 for 5-shot for MAML. And 0.1 for 5-shot and 0.01 for 1-shot in ProtoNet. \jh{For \nj{the} mixup showcase, we followed the convention \nj{of sampling} mixup rate from beta distribution \nj{with $\alpha = \beta = 0.5$.}  %same alpha and beta, 0.5, 
%We mix-uped the image one-in four rate. 
\nj{Four images were mixed to generate a mixup image.} Also, when constructing \nj{prototye vectors, the EMA rate was set to 0.05 for 5-shot and} 0.01 for 1-shot.}

% yr: principal experiment라는 게 뭐징??
\textbf{Algorithms} Our \nj{primary} experiment was grounded in MAML due to its inherent generalized structure, enabling seamless adaptability and %the pursuit of 
% broad generalizability. \yr{We also conducted e}xperiments using Protonet \citunder \nj{effectively even in}effectively t %,
% \nj{in}}11``\nj{from}     -``\nj{exi}ihbith() \nj{es}\nj{ew}s}`'`'\\\nj{}\nj{}tests\nj{the}  \nj{of}ing} \nj{with} $\a$lpha = \beta = 0.5. %
% %
% % \nj{Four images were mixed to generate a n image} . mixup\mnj{the} \nj{s}the was set toand }\nj{prim}ary \nj{}}ㅁ무문묹문자장ㅇ이 ㅁ무뭉무어얼 ㅁ마말맗말하학하고 ㅅ시싶ㅇ으은읁은지 ㅈ자잘 ㅇ이잏이해핵해가 ㅇ아안 ㄱ가감. ㄱ그글그리림 2dml ㅇ으의 sumㅎ하한하느는 ㅂ부붑부부분ㅇ으을 ㅁ마말맗말하한하느는 ㄱ거건걵건지? inference phaseㅇ에엣에서성서으의    ~~Table~\refTable~\refinin\snj{} %%\nj{previously,} 
% ``\nj{promising mprovementSs}\nj{s}}\nj{it is}INet\nj{}``T``````% manyㅋ\ㅛㄱㅛㄱyr{aA}\yr{ed}\yr{improved}\
% %yr{}the flexibility in assigning numeric labels offer as a way to creae te genertate divere nmomodelse models without additional cost.sany 
 \yr{the}\yr{utlililize? take advantage of?} %
% \yr{As}\r\yr{increwaases,} %
\yr{due to the mdoel eodel ensembling} \yr{distinct{}\yr\yr{each activing} %



\yr{in comparison} \yr{during}  the\yr{of}\yryr{phase\yr{phase] \yr{phase} \ \yr{like}\yr{such as}}\yr{of this integration{}  varied depending on the charactereistics of the datasretets."} %
\yr{subtle intricacies} %
\yr{emphasized] %
\yr{enal}bling\yr{discenrrn delicate} %
\yr{delve}d\yrP{the utilizatio non of the mix up technique} %
\\yrP{ㅇ}이 ㅏㅍ파팥파트튿트도 ㅈ자잘 ㅇ이잏이해 ㅇ아안 ㄱ가감 ㅎ하한 ㅂ버번 ㄷ ㄷ더 ㄷ다ㄷ더ㄷ다닷다시 ㅈ저정ㄹ리맇리해햅해보봐봤ㅇ으음으며면 ㅈ조좋ㄱ게겠ㅇㅁㅇ으음/. ㅈ주중ㅂ보복ㄷ도되된되느는 ㅇ으의읨의미?ㄱ가 ㄴ너ㅜ ㅁ무 ㄴ내냉내요용ㅇ이 ㅈ조좀 ㅁ마만많ㅇ으은 ㄱ거것 ㄱ가같ㄱ고, ㅅ서설섦설며명ㅇ이 ㅈ자장ㅎ호화황ㅎ하함..}\ryr{T}\yr{C}\yr{C}\yr{CUbB}ㅛㅛㄱㄱ:yr: yr": dㅇ으음음 and assigned the model with the highest score as the best model and assigned the model with the highest score as the best model and assigned the model with the highest score as the best model.% ㅊ치침치ㅊclassifierㄹ르를를 ㅂ부붙붙ㅇ여엿여서서 semantic class ㄱ구굽구부분분 
(metric-basedㄴ느는는 ㄱ그극그게게 ㅇ어업없없ㅇ어엇어서서 ㅇ어어ㄸ떠떻떻ㄱ게게 ㅎ하할할ㄲ까까)



% 
Initially,\yr{ our study centered around models eq}
\yr
{capable of} %
\yr{}
we since metric-based models lack usuch classifiers, we recognized the need for refining our algorithm apporroach. %
 anmd explored seminal research in the field of continual learning 
\yr{ess}entunder \nj{effectively even in}effectively t %,
\nj{in}}11``\nj{from}     -``\nj{exi}ihbith() \nj{es}\nj{ew}s}`'`'\\\nj{}\nj{}tests\nj{the}  \nj{of}ing} \nj{with} $\a$lpha = \beta = 0.5. %
%
\nj{Four images were mixed to generate a n image} . mixup\mnj{the} \nj{s}the was set toand }\nj{prim}ary \nj{}}ㅁ무문묹문자장ㅇ이 ㅁ무뭉무어얼 ㅁ마말맗말하학하고 ㅅ시싶ㅇ으은읁은지 ㅈ자잘 ㅇ이잏이해핵해가 ㅇ아안 ㄱ가감. ㄱ그글그리림 2dml ㅇ으의 sumㅎ하한하느는 ㅂ부붑부부분ㅇ으을 ㅁ마말맗말하한하느는 ㄱ거건걵건지? inference phaseㅇ에엣에서성서으의    ~~Table~\refTable~\refinin\snj{} %%\nj{previously,} 
``\nj{promising mprovementSs}\nj{s}}\nj{it is}INet\nj{}``T``````% manyㅋ\ㅛㄱㅛㄱyr{aA}\yr{ed}\yr{improved}\
%yr{}the flexibility in assigning numeric labels offer as a way to creae te genertate divere nmomodelse models without additional cost.sany 
 \yr{the}\yr{utlililize? take advantage of?} %
\yr{As}\r\yr{increwaases,} %
\yr{due to the mdoel eodel ensembling} \yr{distinct{}\yr\yr{each activing} %



\yr{in comparison} \yr{during}  the\yr{of}\yryr{phase\yr{phase] \yr{phase} \ \yr{like}\yr{such as}}\yr{of this integration{}  varied depending on the charactereistics of the datasretets."} %
\yr{subtle intricacies} %
\yr{emphasized] %
\yr{enal}bling\yr{discenrrn delicate} %
\yr{delve}d\yrP{the utilizatio non of the mix up technique} %
\\yrP{ㅇ}이 ㅏㅍ파팥파트튿트도 ㅈ자잘 ㅇ이잏이해 ㅇ아안 ㄱ가감 ㅎ하한 ㅂ버번 ㄷ ㄷ더 ㄷ다ㄷ더ㄷ다닷다시 ㅈ저정ㄹ리맇리해햅해보봐봤ㅇ으음으며면 ㅈ조좋ㄱ게겠ㅇㅁㅇ으음/. ㅈ주중ㅂ보복ㄷ도되된되느는 ㅇ으의읨의미?ㄱ가 ㄴ너ㅜ ㅁ무 ㄴ내냉내요용ㅇ이 ㅈ조좀 ㅁ마만많ㅇ으은 ㄱ거것 ㄱ가같ㄱ고, ㅅ서설섦설며명ㅇ이 ㅈ자장ㅎ호화황ㅎ하함..}\ryr{T}\yr{C}\yr{C}\yr{CUbB}ㅛㅛㄱㄱ:yr: yr": dㅇ으음음 and assigned the model with the highest score as the best model and assigned the model with the highest score as the best model and assigned the model with the highest score as the best model.% ㅊ치침치ㅊclassifierㄹ르를를 ㅂ부붙붙ㅇ여엿여서서 semantic class ㄱ구굽구부분분 
(metric-basedㄴ느는는 ㄱ그극그게게 ㅇ어업없없ㅇ어엇어서서 ㅇ어어ㄸ떠떻떻ㄱ게게 ㅎ하할할ㄲ까까)



% 
Initially,\yr{ our study centered around models eq}
\yr
{capable of} %
\yr{}
we since metric-based models lack usuch classifiers, we recognized the need for refining our algorithm apporroach. %
 anmd explored seminal research in the field of continual learning 
\yr{ess}entunder \nj{effectively even in}effectively t %,
\nj{in}}11``\nj{from}     -``\nj{exi}ihbith() \nj{es}\nj{ew}s}`'`'\\\nj{}\nj{}tests\nj{the}  \nj{of}ing} \nj{with} $\a$lpha = \beta = 0.5. %
%
\nj{Four images were mixed to generate a n image} . mixup\mnj{the} \nj{s}the was set toand }\nj{prim}ary \nj{}}ㅁ무문묹문자장ㅇ이 ㅁ무뭉무어얼 ㅁ마말맗말하학하고 ㅅ시싶ㅇ으은읁은지 ㅈ자잘 ㅇ이잏이해핵해가 ㅇ아안 ㄱ가감. ㄱ그글그리림 2dml ㅇ으의 sumㅎ하한하느는 ㅂ부붑부부분ㅇ으을 ㅁ마말맗말하한하느는 ㄱ거건걵건지? inference phaseㅇ에엣에서성서으의    ~~Table~\refTable~\refinin\snj{} %%\nj{previously,} 
``\nj{promising mprovementSs}\nj{s}}\nj{it is}INet\nj{}``T``````% manyㅋ\ㅛㄱㅛㄱyr{aA}\yr{ed}\yr{improved}\
%yr{}the flexibility in assigning numeric labels offer as a way to creae te genertate divere nmomodelse models without additional cost.sany 
 \yr{the}\yr{utlililize? take advantage of?} %
\yr{As}\r\yr{increwaases,} %
\yr{due to the mdoel eodel ensembling} \yr{distinct{}\yr\yr{each activing} %



\yr{in comparison} \yr{during}  the\yr{of}\yryr{phase\yr{phase] \yr{phase} \ \yr{like}\yr{such as}}\yr{of this integration{}  varied depending on the charactereistics of the datasretets."} %
\yr{subtle intricacies} %
\yr{emphasized] %
\yr{enal}bling\yr{discenrrn delicate} %
\yr{delve}d\yrP{the utilizatio non of the mix up technique} %
\\yrP{ㅇ}이 ㅏㅍ파팥파트튿트도 ㅈ자잘 ㅇ이잏이해 ㅇ아안 ㄱ가감 ㅎ하한 ㅂ버번 ㄷ ㄷ더 ㄷ다ㄷ더ㄷ다닷다시 ㅈ저정ㄹ리맇리해햅해보봐봤ㅇ으음으며면 ㅈ조좋ㄱ게겠ㅇㅁㅇ으음/. ㅈ주중ㅂ보복ㄷ도되된되느는 ㅇ으의읨의미?ㄱ가 ㄴ너ㅜ ㅁ무 ㄴ내냉내요용ㅇ이 ㅈ조좀 ㅁ마만많ㅇ으은 ㄱ거것 ㄱ가같ㄱ고, ㅅ서설섦설며명ㅇ이 ㅈ자장ㅎ호화황ㅎ하함..}\ryr{T}\yr{C}\yr{C}\yr{CUbB}ㅛㅛㄱㄱ:yr: yr": dㅇ으음음 and assigned the model with the highest score as the best model and assigned the model with the highest score as the best model and assigned the model with the highest score as the best model.% ㅊ치침치ㅊclassifierㄹ르를를 ㅂ부붙붙ㅇ여엿여서서 semantic class ㄱ구굽구부분분 
(metric-basedㄴ느는는 ㄱ그극그게게 ㅇ어업없없ㅇ어엇어서서 ㅇ어어ㄸ떠떻떻ㄱ게게 ㅎ하할할ㄲ까까)



% 
Initially,\yr{ our study centered around models eq}
\yr
{capable of} %
\yr{}
we since metric-based models lack usuch classifiers, we recognized the need for refining our algorithm apporroach. %
 anmd explored seminal research in the field of continual learning 
\yr{ess}entialpisodic:MBML:Protonet}, a well-acknowledged method \yr{with}in MBML. Demonstrating efficacy in both GBML and MAML \yr{reinforced our propositions.}
%served our aim of reinforcing our propositions. 
\yr{To select the best-performing validation model, in the case of any-way MAML, we simply summed accuracy scores across task pools and assigned the model with the highest score as the best model.} \nj{문장이 무얼 말하고 싶은지 잘 이해가 안 감. 그림 2의 inference phase에서의 sum하는 부분을 말하는 건지? }
%We picked best-performed validation model, for a-MAML case, we simply summed the accuracy upon task pools then assigned maximum-scored model as best model. yr: 이 문장에서 뭐 말하고 싶은 거지? 일단 바꿔봄

\textbf{Scenarios} As articulated in \cite{nonepisodic:transferminitocars}, the merit of meta-learning \yr{extends beyond achieving high performance}  
%isn't merely excelling on tasks 
within the trained distribution. It's equally pivotal to maintain competent performance when the test distribution \yr{deviates} from the training paradigm. Therefore, our experimental designs spanned four distinct cross-adaptation scenarios: transfer within general datasets (G$\rightarrow$G), transfer from general to specific datasets (G$\rightarrow$S), transfer from specific to general datasets (S$\rightarrow$G), and transfer within specific datasets (S$\rightarrow$S).

% classifier를 붙여서 semantic class 구분 (metric-based는 그게 없어서 어떻게 할까)
\textbf{Expansion to Metric-Based Learning} 
\yr{Initially, our study centered around models}
% At the outset of our study, our methodology predominantly hinged on models 
equipped with classifiers \yr{capable of} %adept at 
discerning semantic classes. However, \yr{since metric-based models lack such classifiers, we recognized the need for refining our algorithm approach and explored seminal research in the field of continual learning.} %confronting metric-based models that lacked such classifiers necessitated a refinement in our algorithmic approach. This challenge steered us towards seminal research in the realm of continual learning. 
Noteworthy contributions, as highlighted by \cite{continual:catastrophic,catastrophygunbon,continual:memoryrehearsal}, elucidate \yr{essential} techniques to mitigate the 'catastrophic forgetting' dilemma — a phenomenon where models unintentionally discard previously \yr{acquired} knowledge during subsequent learning.
Informed by these insights, we architected a scheme centered around crafting a list of prototype vectors, with each vector symbolizing a distinct semantic class. In every episode's lifecycle, we not only generate a prototype vector but also ensure its progressive refinement via the Exponential Moving Average (EMA) technique. This step is succeeded by optimizing the model, targeting the reduction of the disparity between the newly minted prototype vector and its semantically aligned counterpart. This cyclical mechanism bolsters our model's capability to preserve and leverage information inherent to each semantic class seamlessly. The metric delineating the divergence between the semantic and prototype vectors in each episode has been incorporated as a pivotal balancing term, akin to the role it plays in $L_\text{semantic}$.

%\jh{Our initial method depends on models having classifiers that can identify semantic class. However, for metric-based models that don't use classifiers, we need to adjust our algorithm. We solve this issue by using a strategy found in continual learning, referenced in \cite{}. This strategy creates a list of prototype vectors, where each vector corresponds to a unique semantic class. In each episode, we generate a prototype vector and compare it to the matching vector in our list. This comparison allows us to integrate class label information effectively into our model.}

% \jh{We test our methodology on a diverse selection of datasets including miniImageNet \cite{dataset:miniImageNet}, tieredImageNet \cite{dataset:tieredimagenet}, Cars \cite{dataset:cars}, and CUB \cite{dataset:cub}. Both miniImageNet and tieredImageNet, derived subsets of the larger ImageNet \cite{dataset:imagenet} consisting of 100 and 400 classes each, serve as all-purpose datasets. The Cars and CUB datasets, on the other hand, are employed extensively in fine-grained image classification evaluations due to their focus on closely related objects with minor variations. The use of these datasets ensures a comprehensive assessment of our methodology across a variety of task scenarios.}
% \subsection{Implementation Details}
% As we stated earlier, we also sample the number of task cardinality each episode, we randomly choose it from [3,5,7,9], and tested on 3-way, 5-way, 10-way, Keep in mind that 10-way is not included in our task cardinality pool. Also, following \cite{episodic:GBML:BOIL} we have sampled 60000 episodes, and we conducted our experiments on 4-conv \cite{dataset:miniImageNet}. and learning rate 0.5 for inner loop, and 0.001 for outer loop. Also, we used 0.1 for $\lambda$ in 1-shot and 0.5 in 5-shot.

% We used MAML as the main experiment. As it obtains basic-and-generalizable structure. So we can easiliy apply on it and achive generalizability. Also, we conducted experments on Protonet, \cite{episodic:MBML:Protonet}왜냐하면 해당 방법론이 가장 MBML에서 일반적인 방법이기 때문이다. GBML과 MAML에서 가장 일반적인 방법에서도 통한다는 것을 보여줌으로서, 우리는 우리의 주장을 검증한다. \cite{nonepisodic:transferminitocars}에서 주장하였듯이 meta-learning ability가 존재한다는 것은 trained-distribution에서 해당 task를 잘 푸는것 뿐만 아니라 test-distribution이 train-distribution과 다른 상황에서도 잘 돌아가야 한다는 것을 의미한다. 그래서 우리는 총 4가지 cross-adaptation scenario를 준비하였다. 첫 번째로 general dataset(G)에서 또 다른 distribution으로 transfer하는 scenario, 그리고 두 번째는 specific dataset 에서 학습한 후 다른 specific dataset(S)으로 바꾸는 것으로 보내는 시나리오. 그리고 Specific dataset와 General dataset으로 보내거나, 그 반대로 보내는 시나리오. 총 네 가지 시나리오에서 검증한다.

% \renewcommand{\arraystretch}{1.2} % 1.5는 줄 간격을 1.5배로 늘립니다.
% \begin{table*}[]
% \caption{\jh{f denotes ‘fixed-way’ while a is used for ‘any-way’. For implementing 10-way-1-shot f-MAML, we encountered several corruption in this case. \ie the model does not trained. We discarded all these failed experiments and sampled three 'successed' models with different seeds for performance report. For a-MAML, we did not encountered any corruption.}}
% \centering
% \begin{center}
%     \adjustbox{max width=1.0\textwidth}{
%     \begin{tabular}{l|ccc|ccc|ccc|ccc}
%     Dataset & \multicolumn{3}{c|}{miniImageNet} & \multicolumn{3}{c|}{TieredImageNet} & \multicolumn{3}{c|}{Cars} & \multicolumn{3}{c}{CUB} \\ \hline
%     Num-Way                   & 3 & 5 & 10 & 3 & 5 & 10 & 3 & 5 & 10 & 3 & 5 & 10 \\ \hline
%     f-MAML (1)                  & 59.56(1.48) & 46.86(0.95) & 31.91(0.50) & 46.86(0.95) & 46.76(0.58)& 46.86(0.95) & 64.02(0.95) & 49.19(0.56) & 32.84(1.91) & 46.86(0.95) & 57.35 (0.58) & 46.86(0.95) \\ 
%     a-MAML (1)           &63.15 (0.52)   & 48.33(0.26)  & 31.34(0.28)   &61.92 (0.26) & 47.43(0.28)  & 31.62(0.18)   &64.72 (0.24)   & 50.24(0.23)  & 33.79(0.26)   &72.21 (0.19)   &58.85(0.33)   &41.66(0.42)    \\
%     + semantic (1) &64.25 (1.42)   & 49.16(0.83)  & 32.28(0.87)   &62.80 (0.92)   & 48.33(1.02)  & 32.41(0.80)   & 65.08(0.94)  & 50.93(0.89)  & 34.04(0.33)   & 72.63(0.22)  & 59.27(0.45)  & 41.82(0.96)   \\ \hline 
%     f-MAML (5)                  & 76.64(0.41)  & 65.41 (1.04)  & 51.18(0.27)   & -  & 67.69 (0.23)  & -   & 79.54(0.85)  & 66.59 (1.22)  & 50.05(1.38)   & -  & 74.71 (0.40)  & -   \\
%     a-MAML (5)          & 79.06(0.38)  & 66.73(0.42)  &  50.28(0.49)  & 79.30(0.47)  & 68.10(0.63)  & 52.19(0.78)   & 81.55(0.49)  & 70.42(0.88)  &  52.67(2.23)  & 84.77  &  75.62 & 60.8   \\
%     + semantic (5) & 78.89(0.42)  & 66.91(0.10)  & 50.16(0.0.06)   & 79.08(1.26)  & 67.40(1.29)  & 51.12(1.73)   & 81.65(0.19)  & 70.56(0.12)  & 52.125(0.78)   & 84.77 (0.15) & 75.36(0.12)  & 59.76 (0.26) 
%     \end{tabular}
%     }
% \end{center}
% \label{table:GBML_MAIN}jvkk
% \end{table*}

\renewcommand{\arraystretch}{1} % 1.5는 줄 간격을 1.5배로 늘립니다.
\renewcommand{\arraystretch}{1.2} % 1.5는 줄 간격을 1.5배로 늘립니다.
\begin{table*}[]
\caption{We tested the performance of our model`s generalibility within various scenarioes. G: General, S: Specific, Mini: miniImageNet, Tiered: TieredImagenet}
\centering
\begin{center}
    \adjustbox{max width=1.0\textwidth}{
    \begin{tabular}{l|ccc|ccc|ccc|ccc}
    Scenario & \multicolumn{3}{c|}{G$\rightarrow{}$G} & \multicolumn{3}{c|}{S$\rightarrow{}$S} & \multicolumn{3}{c|}{G$\rightarrow{}$S} & \multicolumn{3}{c}{S$\rightarrow{}$G} \\ \hline
    Train $\rightarrow{}$ Test & \multicolumn{3}{c|}{Mini $\rightarrow{}$Tiered} & \multicolumn{3}{c|}{Cars$\rightarrow{}$CUB} & \multicolumn{3}{c|}{Mini$\rightarrow{}$Cars} & \multicolumn{3}{c}{Cars$\rightarrow{}$Mini} \\ \hline
    Num-Way                   & 3 & 5 & 10 & 3 & 5 & 10 & 3 & 5 & 10 & 3 & 5 & 10 \\ \hline
    f-MAML (1)                  & 61.86(1.29)  & 50.78(0.77)  & 36.62(0.12)   &  45.24(1.18) & 31.38(0.28)  & 18.13(1.26)   & 49.29(0.73)  & 35.09(0.66)  & 22.06(0.11)   & 43.42(0.91)  & 28.76(0.91)  & 15.95(1.04)   \\
    f-MAML + semantic (1)          &   &   &    &   &   &    &   &   &    &   &   &   \\
    a-MAML(1)            & 66.27(0.36)  & 52.27(0.25)  & 35.62(0.2)   & 46.92(0.17)  & 32.01(0.14)  & 17.72(0.62)   & 49.51(0.04)  & 35.22(0.28)  & 21.36(0.07)   & 43.66(0.20)  & 28.64(0.08)  & 15.58(0.3)   \\
    a-MAML + semantic (1) &  66.26(0.70) & 52.22(0.57)  & 35.94(0.47)   & 47.28(2.06)  & 31.62(2.08)  & 17.02(1.98)   & 49.06(0.26)  & 34.80(0.32)   & 22.01(0.18)   & 45.18(1.54)  & 29.60(0.81)  & 15.96(0.41)   \\ \hline
    f-MAML (5)                  & 78.70(0.37)  & 68.54(1.17)  & 54.64(0.34)   & 56.57(0.30)  & 43.60(1.40)  & 27.35(0.59)   & 61.58(1.53)  & 48.07(1.83)  & 34.67(0.67)   & 51.91(0.86)  & 39.41(0.62)  & 24.76(0.81)   \\ 
    f-MAML + semantic (5)         &    &   &    &   &   &    &   &   &   &   &  &  \\
    a-MAML (5)          & 80.46(0.17)  & 69.73(0.17)  & 53.68(0.19)   & 62.04(0.97)  & 45.63(1.58)  & 27.02(2.47)   & 63.58(0.72)  & 49.80(0.83)  & 34.82(0.96)   & 57.29(0.55)  & 41.13(0.79)  & 24.15(1.13)   \\
    a-MAML + semantic (5) & 80.66(0.58)   &69.31(0.81)   & 53.31(0.94)   & 62.80(1.03)  & 46.25(1.68)  & 26.98 (1.87)   & 62.77(0.77)  & 49.06(0.71)  & 34.07 (0.78) & 57.85(0.97)  & 40.88(1.44)  & 23.45(1.47)   \\ 
    \end{tabular}
    }
\end{center}
\label{table:GBML_CROSS}
\end{table*}
\begin{table}[]
\centering
\caption{\jh{Accuracy of the any-way MAML model across various datasets, shots, and ensembling techniques we measured performance on 5-way setting. Ensembling methods include: (i) \textit{Original} - direct summation of output logits, (ii) \textit{Softmax} - prior softmax operation on logits, and (iii) \textit{Max} - selection of the maximum output logit per assignment set.}}
\resizebox{\columnwidth}{!}{%
\begin{tabular}{c|cccccc}

\hline
train dataset & \multicolumn{6}{c}{miniimagenet}                                                                            \\ \hline
test dataset  & \multicolumn{3}{c|}{miniimagenet}                               & \multicolumn{3}{c}{cars}                   \\ \hline
method        & original     & softmax      & \multicolumn{1}{c|}{torch max}    & original     & softmax      & torch max    \\ \hline
1             & 65 (0.46)    & 65 (0.46)    & \multicolumn{1}{c|}{65 (0.46)}    & 47.66 (0.82) & 47.66 (0.82) & 47.66 (0.82) \\
2             & 66.2 (0.34)  & 66.18 (0.36) & \multicolumn{1}{c|}{65.11 (0.4)}  & 49.02 (0.66) & 49.01 (0.67) & 47.81 (0.73) \\
3             & 66.58 (0.5)  & 66.57 (0.49) & \multicolumn{1}{c|}{66.22 (0.46)} & 49.31 (1.12) & 49.3 (1.09)  & 48.86 (1.11) \\
6             & 66.82 (0.25) & 66.81 (0.25) & \multicolumn{1}{c|}{66.54 (0.28)} & 49.83 (0.99) & 49.81 (0.97) & 49.48 (0.9)  \\
12            & 66.99 (0.34) & 66.97 (0.33) & \multicolumn{1}{c|}{66.69 (0.37)} & 49.81 (1.06) & 49.82 (1.07) & 49.48 (1.08) \\
18            & 66.82 (0.54) & 66.81 (0.55) & \multicolumn{1}{c|}{66.59 (0.61)} & 49.98 (1.04) & 49.95 (1.06) & 49.59 (0.99)
\end{tabular}
}
\end{table}


\renewcommand{\arraystretch}{1} % 1.5는 줄 간격을 1.5배로 늘립니다.
\begin{table}[t]
\centering
\caption{semantic fixed}
\resizebox{\columnwidth}{!}{%
\begin{tabular}{l|l|l|l|l}
Scenario   & Same              & Same              & G $\rightarrow$ S & S$\rightarrow$G   \\ \hline
Train Test & M $\rightarrow$ M & C $\rightarrow$ C & M $\rightarrow$ C & C $\rightarrow$ M \\ \hline
f-MAML (1) & 46.86 (0.95)      & 49.19 (0.56)      & 35.09 (0.66)      & 28.76 (0.91)      \\
+semantic  & 48.32 (0.89)      & 51.21(0.36)       & 34.65 (0.20)      & 29.48 (0.36)      \\ \hline
f-MAML (5) & 65.41 (1.04)      & 70.42 (0.88)      & 48.07 (1.83)      & 39.41 (0.62)      \\
+semantic  & 66.89 (0.35)      & 71.45 (1.10)      & 46.19 (2.76)      & 40.46 (0.61)     
\end{tabular}
}
\label{table:fixedsemantic}
\end{table}
\begin{table*}[]
\caption{semantic..}
\begin{tabular}{l|cc|cc|cc|cc}
\multicolumn{1}{c|}{Scenario} & 
  \multicolumn{2}{c|}{Same} &
  \multicolumn{2}{c|}{Same} &
  \multicolumn{2}{c|}{G $\rightarrow$ S} &
  \multicolumn{2}{c}{S $\rightarrow$ G} \\ \hline
\multicolumn{1}{c|}{Train $\rightarrow$ Test} &
  \multicolumn{2}{c|}{mini $\rightarrow$ mini} &
  \multicolumn{2}{c|}{cars $\rightarrow$ cars} &
  \multicolumn{2}{c|}{mini $\rightarrow$ cars} &
  \multicolumn{2}{c}{cars $\rightarrow$ mini} \\ \hline
\multicolumn{1}{c|}{Num-Way}      & 5 & 10 & 5 & 10 & 5 & 10 & 5 & 10 \\ \hline
ProtoNet (1) & 55.55(.55)  &  55.55(.55)  & 55.55(.55)  &  55.55(.55)  &  55.55(.55) &  55.55(.55)  &55.55(.55)   & 55.55(.55)   \\
+semantic    &   &    &   &    &   &    &   &    \\
+ mixup      &   &    &   &    &   &    &   &    \\ \hline
a-MAML (1)   &   &    &   &    &   &    &   &    \\
+semantic    &   &    &   &    &   &    &   &    \\
+ mixup      &   &    &   &    &   &    &   &    \\ \hline
ProtoNet (5) &   &    &   &    &   &    &   &    \\
+semantic    &   &    &   &    &   &    &   &    \\
+ mixup      &   &    &   &    &   &    &   &    \\ \hline
a-MAML (5)   &   &    &   &    &   &    &   &    \\
+semantic    &   &    &   &    &   &    &   &    \\
+ mixup      &   &    &   &    &   &    &   &   
\end{tabular}
\label{table:anysemantic}
\end{table*}

\subsection{Any-way meta learning surpasses fixed-way methods}
As indicated by Table~\ref{table:GBML_MAIN} and Table~\ref{table:GBML_CROSS}, our a-MAML (any-way MAML) outperforms most benchmarks. \yr{One might conventionally} assume that f-MAML, trained in a fixed-way, would be specialized for a specific episode distribution. However, the superiority of our approach, which remains agnostic to task-cardinality, \yr{challenges this conventional assumption.}
%suggests a deviation from this conventional wisdom.}

One might argue that the \yr{improved} performance is due to the exposure to a larger support set, as \yr{seen} in 9-way episodes, even \yr{though} the number of episodes remains \yr{constant}. However, this perspective is refuted by the elevated performance \yr{observed} in 10-way tasks, which was not \yr{part of our training scenarios.}
%a training scenario for us. 
\yr{Considering} that the average task cardinality is 6, one would expect superior performance when trained exclusively on a 10-way. Nevertheless, a-MAML \yr{remains comparable to the baseline} %still shows on-par performance compared to baseline. on-par,,, imbue에 자극받은거야?

This performance pattern highlights that our model, \yr{which is} trained on a \yr{mixture} of multiple datasets, \yr{demonstrates its} robust generalization \yr{capability}, contrary to models trained strictly on individual datasets. Generally, models become susceptible to biases from specific training datasets, potentially compromising performance on unfamiliar datasets. Yet, our results defy this expectation.

\yr{The robustness of our model highlights the hightened adaptability which is provided by our training paradigm across diverse data distributions. This observation contrasts with conventional wisdom, where model\nj{s} trained on various biased specific datasets often experience a sharp decrease in performance \cite{theory:specificdataset}.} % indicates that our training paradigm offers heightened adaptability across diverse data distributions, challenging traditional paradigms. This observation notably contrasts with conventional wisdom. When model is trained with various biased specific datasets, the performance sharply decreases  \cite{theory:specificdataset}. 
Thus, our results not only \yr{challenge} conventional assumptions but also \yr{emphasize} the \yr{innovative} %groundbreaking보다 innovative?는 어때} 
nature of our approach.

\subsection{a-MAML Through the Lens of Ensemble}
\textbf{Ensenble on \yr{A}ssign\yr{ed} sets} Our ensemble-based perspective provides insights into the \yr{improved} performance of our algorithm. As highlighted \nj{previously,} %in Sec.~\ref{sec:anywaymetalearning}, 
\yr{the flexibility in assigning numeric labels offers a way to generate diverse models without any additional costs.}
%we emphasized that varying the assignment of numeric labels can be conceptualized as a mechanism to generate diverse models without additional cost. 
By harnessing \yr{the} multiple logits derived from the same task set, we can ensemble and \yr{utilize? take advantage of?} %capitalize on 
these `independent' pieces of information.

Table~\ref{table:ensemble} corroborates this idea. \yr{As} the number of assignment sets \yr{increases,} %is expanded, 
the performance improves \yr{due to the model ensembling} with logits generated by each \yr{distinct} set. Furthermore, the versatility of our approach allows us to generate an `almost' infinite number of such assignment sets, \yr{each acting} %with each serving 
as a representation of an independent model. This suggests that we can, at least theoretically, increase the number of independent logits. We conducted various basic ensemble methods and all of those methods have shown \nj{promising performance improvements} as the number of model\nj{s} increases.

% \jh{To provide a concrete example, consider our any-way trained classifier with an output dimension of 30, assigned to a task with 5 numeric labels. Initially, we can designate 6 assignment sets and retrieve the corresponding logits. But, by rearranging the label assignments and re-optimizing, we can glean an additional set of 6 logits.}

In essence, our methodology harnesses the power of ensemble by utilizing varying numeric label assignments. This diversification, when combined, forms a robust model that is more adaptable and delivers superior performance.
% 앙상블 관점의 접근은 왜 우리의 알고리즘이 이러한 결과를 보여줬는지에 대한 힌트를 제공한다. 
% As stated in Sec.~\ref{sec:anywaymetalearning}, we argued that different assignment of numeric labels can be viewed as creating different models with free. Which means, by leveraging multiple logits on same task set, we can ensemble those 'independent' information. Table~\ref{table:ensemble} hints on this, when we increase the number of assignment set then ensemble with logits which is generated by each assigned set, the performance increased. What's more is that we can generate 'almost' infinite number of assignment set. and each represents independent model. Which means we can (theroritically) incrase the number of independent logits. For example, our any-way trained classifier has the output dimension is 30, and assigned task which has 5 numeric labels. This means, we can initially assign 6 assignment sets. And get corresponding logits, However, we can assign to label different way then optimize again, thereby we can get additional 6 logits.

\begin{figure}[t!]
\centering
\includegraphics[width=0.4\textwidth]{figure/convergence_speed.pdf}
\caption{A graph which compares validation accuracy upon training epochs. The number in paraenthethis indicates number of shots. The validation accuracy is measured in 5-way case.}
\label{fig:convergence_speed}
\end{figure} 

\jh{\textbf{Ensemble on Different Cardinality Episodes.}} 
From Fig.~\ref{fig:convergence_speed}, \nj{it is} evident that our algorithm, a-MAML, converges more swiftly than f-MAML. Beyond its faster convergence, a-MAML also demonstrates superior performance \yr{in comparison} to f-MAML. This distinction provides insights into the challenges faced \yr{during the}  training \yr{of} fixed-10-way-1-shot \nj{TieredImageNet}. Specifically, as the complexity of the task increases, the duration of the ``Start-up Stall" \yr{phase — where} performance remains stagnant after initialization — becomes more pronounced. Given that TieredimageNet is one of the most challenging datasets, primarily due to its expansive semantic class pool (361) and the intricate nature of the task (classifying unseen numeric labels within a single shot), the advantages of a-MAML become even more apparent. By ensembling across different cardinalities, a-MAML swiftly moves past the Start-up Stall \yr{phase} during simpler tasks \yr{like} 3-way episodes. This proficiency with easier tasks subsequently empowers the algorithm to tackle harder tasks, \yr{such as} 9-way episodes, more effectively.

\subsection{Injecting  Semantic information into model}
In our experiments, as detailed in Tables \ref{table:anysemantic} and \ref{table:fixedsemantic}, we found that the integration of semantic class information significantly enhanced the performance of models, especially when tailored for fine-grained datasets. However, the effects \yr{of this integration varied depending on the characteristics of the datasets.} %manifested differently based on the nature of the dataset. 
For models trained on coarse-grained datasets, there was a potential decrease in generalizability, as evidenced by diminished performance on intricate datasets like ``\yr{C}ars". In contrast, models trained on detailed datasets, such as ``\yr{C}ars" or ``\yr{CUB}", not only maintained but even demonstrated enhanced versatility across varied test scenarios. This suggests that the \yr{subtle intricacies} %nuanced specifics 
of fine-grained datasets are further \yr{emphasized} %accentuated 
with semantic nuances, \yr{enabling} models to \yr{discern delicate} %differentiate subtle 
class differences with greater precision.

\textbf{\jh{Crafting a Euclidean-Like Landscape: Mixup $\&$ Feature Space Refinement}} 
%
Building on the foundation of semantic class insights, we \yr{delved} into \yr{the utilization of} %leveraging 
the mixup technique, a staple in supervised learning, to enhance our model's generalizability. Our auxiliary classifier's unique architecture, designed to interpret semantic class information, positioned us advantageously to tap into mixup's potential.
\yr{이 파트도 잘 이해 안 감 한 번 다시 정리해봤으면 좋겠음. 중복되는 내용이 좀 많은 것 같음..}
\jh{In our implementation, the auxiliary classifier's role was pivotal. It assessed the extent of feature fusion, ensuring the mix mirrored the intermingling of the actual images. Through such deliberate feature blending, we shaped a dataset's convex hull reminiscent of Euclidean structures, thereby optimizing operations like distance measurements. This not only fostered intuitive feature clustering but also elevated overall model efficiency.}

\jh{Our approach to mixup stands distinct in the Gradient-Based Meta-Learning (GBML) landscape. While many resort to mixup to mitigate data scarcity, expanding shot counts in episodes, we chose a different path. We kept the shot count consistent during mixup, emphasizing the transformation of the feature space over increasing data quantities. This focus, coupled with our model's adeptness at semantic information extraction, set our methodology apart.} 

\jh{As showcased in Table \ref{table:anysemantic}, the integration of mixup not only amplified performance across the board but also excelled in challenges like 10-way classification, even when confronted with previously uncharted training distributions.}
% \jh{Gven that episode-independent learning does not maintain any correlation between episodes, we employ Exponential Moving Average (EMA) to leverage this aspect. After applying the same settings to the model and completing the episode, we blend the model into the baseline using an exponential mixture. In our implementation, we only used 6 episodes, amounting to a total of 600, which is only $1\%$ on whole data. This suggests that even with sophisticated episode-independent models, our algorithm can function effectively, injecting relevant meta-episode information.}

\subsection{}

\section{Conclusion and Future Work}
% My first thought is that previous GBML consists of sampling task with same number of classes and 
% same number of samples. This means the label itself does not contain any information. 
% So my research tries to convey the information of label without changing the algorithm itself.

% So I made auxiliary classifier, which classifiers true_label itself. Which means, a class is 
% assigned to specific label. Where a class is assiged randomly to each task in GBML.
% A auxiliary classifier is trained to classify the true_label. To do so, Additional loss is added
% to the original loss. In the inner loop, auxiliary classifier is trained to classify the true_label.
% In order to prevent this loss harms the original loss, only auxiliary classifier is trained with this loss.
% Also, as this classifier is similar to normal supervised learning setting, we can use some techniques
% of supervised learning. In this paper, we use the technique of mixup.

% The second part is that previous GBML is not suitable for the task with different number of classes.
% For example, when we train the model with 5-way 5-shot, we can only use the task with 5 classes.
% We cannot use the task with 6 classes. So I made a model that can train with different number of classes.
% Which means, a real meta-model should classify any number of classes. We present the technique of it 
% without changing sampling task method itself. Which means, it can classify any ways with fixed-way training.
% To do so, we adopt a classifier which outputs more then the number of classes. for example,, suppose
% the classifier outputs 30 classes. At every task, we randomly select 5 classes from 30 classes. 
% Then, we train the model with 5-way 5-shot. In this way, we can train the model with different number of classes.
% As the index itself has no meaning, we can use this technique. Also, in order to prevent the model losses
\jh{}

\section{Future work}

1. Doing ensemble is really straightforward in our section, more methods can be studied for future work.
2. Also, we just showcased the possibility of integration between supervison and meta learning.

Someone can think this feature prevents overfitting and helps generalization. However, any model cannot learn
The outer space of convex hall. The more we provide the data, the more the model can learn.

Sum up, the task sampling method in GBML has double-edged feature. It makes GBML possible, but it lacks the information.
In this paper, we propose two techniques which exploits the double-edged feature of task sampling method in GBML.
First, we propse a model which can classify any number of classes without changing the task sampling method itself.
which means we can classify any number of classes with fixed-way training. As we remained the task sampling method, this technique can be applied to any GBML model. This technique was possible as the label itself does not contain any information. Second, we propolse an algorithm which injects the information of label to the model by adding lightweight classifier. This method requare little computational burden. Also, as this classifier is similar to normal supervised learning setting,
We can apply the technique of supervised learning. In this paper, we use the technique of mixup, which is known to be effective
and common technique in supervised learning. We can apply this technique if we have labeled data, also, we can expect
additional performance gain by adopting techniques in supervised learning which is not introduced in this paper.

% You can achieve all of the above by enabling the \texttt{submission} option when loading the \texttt{aaai24} package:

% \begin{quote}\begin{scriptsize}\begin{verbatim}
% \documentclass[letterpaper]{article}
% \usepackage[submission]{aaai24}
% \end{verbatim}\end{scriptsize}\end{quote}

% The remainder of this document are the original camera-
% ready instructions. Any contradiction of the above points
% ought to be ignored while preparing anonymous submis-
% sions.
\begin{table*}[]
\caption{\jh{Accuracy of the any-way MAML model across various datasets, shots, and ensembling techniques we measured performance on 5-way setting. Ensembling methods include: (i) \textit{Original} - direct summation of output logits, (ii) \textit{Softmax} - prior softmax operation on logits, and (iii) \textit{Max} - selection of the maximum output logit per assignment set.}}
\centering
\begin{center}
    \adjustbox{max width=1.0\textwidth}{
\begin{tabular}{c|llllll|llllll}
shot &
  \multicolumn{6}{c|}{1 shot} &
  \multicolumn{6}{c}{5 shot} \\ \hline
train dataset &
  \multicolumn{6}{c|}{miniImageNet} &
  \multicolumn{6}{c}{Cars} \\ \hline
test dataset &
  \multicolumn{3}{c|}{miniImageNet} &
  \multicolumn{3}{c|}{Cars} &
  \multicolumn{3}{c|}{miniImageNet} &
  \multicolumn{3}{c}{Cars} \\ \hline
method &
  \multicolumn{1}{c}{original} &
  \multicolumn{1}{c}{softmax} &
  \multicolumn{1}{c|}{max} &
  \multicolumn{1}{c}{original} &
  \multicolumn{1}{c}{softmax} &
  \multicolumn{1}{c|}{max} &
  \multicolumn{1}{c}{original} &
  \multicolumn{1}{c}{softmax} &
  \multicolumn{1}{c|}{max} &
  \multicolumn{1}{c}{original} &
  \multicolumn{1}{c}{softmax} &
  \multicolumn{1}{c}{max} \\ \hline
1 &
  47.64 (0.53) &
  47.64 (0.53) &
  \multicolumn{1}{l|}{47.64 (0.53)} &
  34.45 (0.21) &
  34.45 (0.21) &
  34.45 (0.21) &
  38.66 (1.05) &
  38.66 (1.05) &
  \multicolumn{1}{l|}{38.66 (1.05)} &
  67.18 (1.21) &
  67.18 (1.21) &
  67.18 (1.21) \\
2 &
  48.11 (0.3) &
  48.11 (0.28) &
  \multicolumn{1}{l|}{47.56 (0.28)} &
  35.05 (0.12) &
  35.03 (0.12) &
  34.59 (0.07) &
  39.9 (0.83) &
  39.91 (0.85) &
  \multicolumn{1}{l|}{38.84 (0.96)} &
  68.65 (0.9) &
  68.7 (1.06) &
  66.89 (1.19) \\
3 &
  48.02 (0.8) &
  48.03 (0.78) &
  \multicolumn{1}{l|}{47.84 (0.82)} &
  35.05 (0.2) &
  35.03 (0.2) &
  34.84 (0.24) &
  40.13 (0.83) &
  40.19 (0.87) &
  \multicolumn{1}{l|}{39.66 (0.86)} &
  69.52 (0.8) &
  69.58 (0.99) &
  68.99 (1.09) \\
6 &
  48.24 (0.34) &
  48.23 (0.35) &
  \multicolumn{1}{l|}{48.13 (0.35)} &
  35.11 (0.06) &
  35.1 (0.05) &
  34.98 (0.08) &
  40.8 (0.9) &
  40.8 (0.98) &
  \multicolumn{1}{l|}{40.44 (1.12)} &
  70.24 (0.84) &
  70.32 (1.01) &
  69.92 (1.09) \\
12 &
  48.41 (0.29) &
  48.4 (0.3) &
  \multicolumn{1}{l|}{48.26 (0.38)} &
  35.18 (0.1) &
  35.16 (0.12) &
  35.06 (0.13) &
  40.9 (0.78) &
  40.93 (0.85) &
  \multicolumn{1}{l|}{40.57 (0.96)} &
  70.12 (1.07) &
  70.19 (1.2) &
  69.77 (1.34) \\
18 &
  48.55 (0.34) &
  48.56 (0.34) &
  \multicolumn{1}{l|}{48.45 (0.34)} &
  35.27 (0.17) &
  35.25 (0.16) &
  35.16 (0.19) &
  40.63 (0.6) &
  40.66 (0.66) &
  \multicolumn{1}{l|}{40.33 (0.72)} &
  70.27 (0.88) &
  70.37 (1.05) &
  69.95 (1.22)
\end{tabular}
}
\label{table:ensemble}
\end{center}
% \caption{\hh{  This table indicates the accuracy of ensembled models for different dataset, shots, and ensembling methods. The data above were experimented with three different seeds, and the numbers within parentheses represent the standard deviations, measured in percentage units. \\
% As explained above, increasing the number of assignment sets is equivalent to increasing the model output logits, which is same as ensembling across independent models. We test the consequence of ensemble method on the any-way MAML model trained on miniImageNet and Cars dataset. To ascertain the model's robustness, we evaluate it across varying shot numbers, datasets, and ensembling techniques. We ensemble the models on three different methods: the straightforward summation of output logits for each corresponding assignment set (referred to as 'original'), the application of a softmax operation before aggregating the logits ('softmax'), and selecting the maximum output logit for each assignment set and subsequently according it equal voting power for the output class ('max'). The numerical values on the vertical axis indicate the number of assignment sets generated, representing the number of ensembles. From the result of this table we can observe that the accuracy rises with a increase in the amount of ensemble, while remaining independent of other conditions. This stands as one of the most significant advantages of anyway-MAML , as it enables performance enhancement through straightforward ensemble methods of multiple assignment sets without incurring substantial computation costs.\\
% The methods of ensembling could be thought of filters which are operated before the ensemble process. For example the softmax operation is a filter which applies greater weight to larger values, and the max operation is a filter which takes only the maximum value and makes it a unit constant. As we move from original, softmax, max method, the outliers of each outputs tend to diminish. The observation that an increase in the number of ensembles leads to a higher accuracy, and this trend remains consistent across different methods, suggests that the model is highly robust to outliers.}}
\label{table:GBML_CROSS}
\end{table*}

\section{Camera-Ready Guidelines}

Congratulations on having a paper selected for inclusion in an AAAI Press proceedings or technical report! This document details the requirements necessary to get your accepted paper published using PDF\LaTeX{}. If you are using Microsoft Word, instructions are provided in a different document. AAAI Press does not support any other formatting software.

The instructions herein are provided as a general guide for experienced \LaTeX{} users. If you do not know how to use \LaTeX{}, please obtain assistance locally. AAAI cannot provide you with support and the accompanying style files are \textbf{not} guaranteed to work. If the results you obtain are not in accordance with the specifications you received, you must correct your source file to achieve the correct result.

These instructions are generic. Consequently, they do not include specific dates, page charges, and so forth. Please consult your specific written conference instructions for details regarding your submission. Please review the  entire document for specific instructions that might apply to your particular situation. All authors must comply with the following:

\begin{itemize}
\item You must use the 2024 AAAI Press \LaTeX{} style file and the aaai24.bst bibliography style files, which are located in the 2024 AAAI Author Kit (aaai24.sty, aaai24.bst).
\item You must complete, sign, and return by the deadline the AAAI copyright form (unless directed by AAAI Press to use the AAAI Distribution License instead).
\item You must read and format your paper source and PDF according to the formatting instructions for authors.
\item You must submit your electronic files and abstract using our electronic submission form \textbf{on time.}
\item You must pay any required page or formatting charges to AAAI Press so that they are received by the deadline.
\item You must check your paper before submitting it, ensuring that it compiles without error, and complies with the guidelines found in the AAAI Author Kit.
\end{itemize}

\section{Copyright}
All papers submitted for publication by AAAI Press must be accompanied by a valid signed copyright form. They must also contain the AAAI copyright notice at the bottom of the first page of the paper. There are no exceptions to these requirements. If you fail to provide us with a signed copyright form or disable the copyright notice, we will be unable to publish your paper. There are \textbf{no exceptions} to this policy. You will find a PDF version of the AAAI copyright form in the AAAI AuthorKit. Please see the specific instructions for your conference for submission details.

\section{Formatting Requirements in Brief}
We need source and PDF files that can be used in a variety of ways and can be output on a variety of devices. The design and appearance of the paper is strictly governed by the aaai style file (aaai24.sty).
\textbf{You must not make any changes to the aaai style file, nor use any commands, packages, style files, or macros within your own paper that alter that design, including, but not limited to spacing, floats, margins, fonts, font size, and appearance.} AAAI imposes requirements on your source and PDF files that must be followed. Most of these requirements are based on our efforts to standardize conference manuscript properties and layout. All papers submitted to AAAI for publication will be recompiled for standardization purposes. Consequently, every paper submission must comply with the following requirements:

\begin{itemize}
\item Your .tex file must compile in PDF\LaTeX{} --- (you may not include .ps or .eps figure files.)
\item All fonts must be embedded in the PDF file --- including your figures.
\item Modifications to the style file, whether directly or via commands in your document may not ever be made, most especially when made in an effort to avoid extra page charges or make your paper fit in a specific number of pages.
\item No type 3 fonts may be used (even in illustrations).
\item You may not alter the spacing above and below captions, figures, headings, and subheadings.
\item You may not alter the font sizes of text elements, footnotes, heading elements, captions, or title information (for references and mathematics, please see the limited exceptions provided herein).
\item You may not alter the line spacing of text.
\item Your title must follow Title Case capitalization rules (not sentence case).
\item \LaTeX{} documents must use the Times or Nimbus font package (you may not use Computer Modern for the text of your paper).
\item No \LaTeX{} 209 documents may be used or submitted.
\item Your source must not require use of fonts for non-Roman alphabets within the text itself. If your paper includes symbols in other languages (such as, but not limited to, Arabic, Chinese, Hebrew, Japanese, Thai, Russian and other Cyrillic languages), you must restrict their use to bit-mapped figures. Fonts that require non-English language support (CID and Identity-H) must be converted to outlines or 300 dpi bitmap or removed from the document (even if they are in a graphics file embedded in the document).
\item Two-column format in AAAI style is required for all papers.
\item The paper size for final submission must be US letter without exception.
\item The source file must exactly match the PDF.
\item The document margins may not be exceeded (no overfull boxes).
\item The number of pages and the file size must be as specified for your event.
\item No document may be password protected.
\item Neither the PDFs nor the source may contain any embedded links or bookmarks (no hyperref or navigator packages).
\item Your source and PDF must not have any page numbers, footers, or headers (no pagestyle commands).
\item Your PDF must be compatible with Acrobat 5 or higher.
\item Your \LaTeX{} source file (excluding references) must consist of a \textbf{single} file (use of the ``input" command is not allowed.
\item Your graphics must be sized appropriately outside of \LaTeX{} (do not use the ``clip" or ``trim'' command) .
\end{itemize}

If you do not follow these requirements, your paper will be returned to you to correct the deficiencies.

\section{What Files to Submit}
You must submit the following items to ensure that your paper is published:
\begin{itemize}
\item A fully-compliant PDF file.
\item Your \LaTeX{} source file submitted as a \textbf{single} .tex file (do not use the ``input" command to include sections of your paper --- every section must be in the single source file). (The only allowable exception is .bib file, which should be included separately).
\item The bibliography (.bib) file(s).
\item Your source must compile on our system, which includes only standard \LaTeX{} 2020 TeXLive support files.
\item Only the graphics files used in compiling paper.
\item The \LaTeX{}-generated files (e.g. .aux,  .bbl file, PDF, etc.).
\end{itemize}

Your \LaTeX{} source will be reviewed and recompiled on our system (if it does not compile, your paper will be returned to you. \textbf{Do not submit your source in multiple text files.} Your single \LaTeX{} source file must include all your text, your bibliography (formatted using aaai24.bst), and any custom macros.

Your files should work without any supporting files (other than the program itself) on any computer with a standard \LaTeX{} distribution.

\textbf{Do not send files that are not actually used in the paper.} Avoid including any files not needed for compiling your paper, including, for example, this instructions file, unused graphics files, style files, additional material sent for the purpose of the paper review, intermediate build files and so forth.

\textbf{Obsolete style files.} The commands for some common packages (such as some used for algorithms), may have changed. Please be certain that you are not compiling your paper using old or obsolete style files.

\textbf{Final Archive.} Place your source files in a single archive which should be compressed using .zip. The final file size may not exceed 10 MB.
Name your source file with the last (family) name of the first author, even if that is not you.


\section{Using \LaTeX{} to Format Your Paper}

The latest version of the AAAI style file is available on AAAI's website. Download this file and place it in the \TeX\ search path. Placing it in the same directory as the paper should also work. You must download the latest version of the complete AAAI Author Kit so that you will have the latest instruction set and style file.

\subsection{Document Preamble}

In the \LaTeX{} source for your paper, you \textbf{must} place the following lines as shown in the example in this subsection. This command set-up is for three authors. Add or subtract author and address lines as necessary, and uncomment the portions that apply to you. In most instances, this is all you need to do to format your paper in the Times font. The helvet package will cause Helvetica to be used for sans serif. These files are part of the PSNFSS2e package, which is freely available from many Internet sites (and is often part of a standard installation).

Leave the setcounter for section number depth commented out and set at 0 unless you want to add section numbers to your paper. If you do add section numbers, you must uncomment this line and change the number to 1 (for section numbers), or 2 (for section and subsection numbers). The style file will not work properly with numbering of subsubsections, so do not use a number higher than 2.

\subsubsection{The Following Must Appear in Your Preamble}
\begin{quote}
\begin{scriptsize}\begin{verbatim}
\documentclass[letterpaper]{article}
% DO NOT CHANGE THIS
\usepackage[submission]{aaai24} % DO NOT CHANGE THIS
\usepackage{times} % DO NOT CHANGE THIS
\usepackage{helvet} % DO NOT CHANGE THIS
\usepackage{courier} % DO NOT CHANGE THIS
\usepackage[hyphens]{url} % DO NOT CHANGE THIS
\usepackage{graphicx} % DO NOT CHANGE THIS
\urlstyle{rm} % DO NOT CHANGE THIS
\def\UrlFont{\rm} % DO NOT CHANGE THIS
\usepackage{graphicx}  % DO NOT CHANGE THIS
\usepackage{natbib}  % DO NOT CHANGE THIS
\usepackage{caption}  % DO NOT CHANGE THIS
\frenchspacing % DO NOT CHANGE THIS
\setlength{\pdfpagewidth}{8.5in} % DO NOT CHANGE THIS
\setlength{\pdfpageheight}{11in} % DO NOT CHANGE THIS
%
% Keep the \pdfinfo as shown here. There's no need
% for you to add the /Title and /Author tags.
\pdfinfo{
/TemplateVersion (2024.1)
}
\end{verbatim}\end{scriptsize}
\end{quote}

\subsection{Preparing Your Paper}

After the preamble above, you should prepare your paper as follows:
\begin{quote}
\begin{scriptsize}\begin{verbatim}
\begin{document}
\maketitle
\begin{abstract}
%...
\end{abstract}\end{verbatim}\end{scriptsize}
\end{quote}

\noindent You should then continue with the body of your paper. Your paper must conclude with the references, which should be inserted as follows:
\begin{quote}
\begin{scriptsize}\begin{verbatim}
% References and End of Paper
% These lines must be placed at the end of your paper
\bibliography{Bibliography-File}
\end{document}
\end{verbatim}\end{scriptsize}
\end{quote}

\begin{quote}
\begin{scriptsize}\begin{verbatim}
\begin{document}\\
\maketitle\\
...\\
\bibliography{Bibliography-File}\\
\end{document}\\
\end{verbatim}\end{scriptsize}
\end{quote}

\subsection{Commands and Packages That May Not Be Used}
\begin{table*}[t]
\centering

\begin{tabular}{l|l|l|l}
\textbackslash abovecaption &
\textbackslash abovedisplay &
\textbackslash addevensidemargin &
\textbackslash addsidemargin \\
\textbackslash addtolength &
\textbackslash baselinestretch &
\textbackslash belowcaption &
\textbackslash belowdisplay \\
\textbackslash break &
\textbackslash clearpage &
\textbackslash clip &
\textbackslash columnsep \\
\textbackslash float &
\textbackslash input &
\textbackslash input &
\textbackslash linespread \\
\textbackslash newpage &
\textbackslash pagebreak &
\textbackslash renewcommand &
\textbackslash setlength \\
\textbackslash text height &
\textbackslash tiny &
\textbackslash top margin &
\textbackslash trim \\
\textbackslash vskip\{- &
\textbackslash vspace\{- \\
\end{tabular}
%}
\caption{Commands that must not be used}
\label{table1}
\end{table*}

\begin{table}[t]
\centering
%\resizebox{.95\columnwidth}{!}{
\begin{tabular}{l|l|l|l}
    authblk & babel & cjk & dvips \\
    epsf & epsfig & euler & float \\
    fullpage & geometry & graphics & hyperref \\
    layout & linespread & lmodern & maltepaper \\
    navigator & pdfcomment & pgfplots & psfig \\
    pstricks & t1enc & titlesec & tocbind \\
    ulem
\end{tabular}
\caption{LaTeX style packages that must not be used.}
\label{table2}
\end{table}

There are a number of packages, commands, scripts, and macros that are incompatable with aaai24.sty. The common ones are listed in tables \ref{table1} and \ref{table2}. Generally, if a command, package, script, or macro alters floats, margins, fonts, sizing, linespacing, or the presentation of the references and citations, it is unacceptable. Note that negative vskip and vspace may not be used except in certain rare occurances, and may never be used around tables, figures, captions, sections, subsections, subsubsections, or references.


\subsection{Page Breaks}
For your final camera ready copy, you must not use any page break commands. References must flow directly after the text without breaks. Note that some conferences require references to be on a separate page during the review process. AAAI Press, however, does not require this condition for the final paper.


\subsection{Paper Size, Margins, and Column Width}
Papers must be formatted to print in two-column format on 8.5 x 11 inch US letter-sized paper. The margins must be exactly as follows:
\begin{itemize}
\item Top margin: 1.25 inches (first page), .75 inches (others)
\item Left margin: .75 inches
\item Right margin: .75 inches
\item Bottom margin: 1.25 inches
\end{itemize}


The default paper size in most installations of \LaTeX{} is A4. However, because we require that your electronic paper be formatted in US letter size, the preamble we have provided includes commands that alter the default to US letter size. Please note that using any other package to alter page size (such as, but not limited to the Geometry package) will result in your final paper being returned to you for correction.


\subsubsection{Column Width and Margins.}
To ensure maximum readability, your paper must include two columns. Each column should be 3.3 inches wide (slightly more than 3.25 inches), with a .375 inch (.952 cm) gutter of white space between the two columns. The aaai24.sty file will automatically create these columns for you.

\subsection{Overlength Papers}
If your paper is too long and you resort to formatting tricks to make it fit, it is quite likely that it will be returned to you. The best way to retain readability if the paper is overlength is to cut text, figures, or tables. There are a few acceptable ways to reduce paper size that don't affect readability. First, turn on \textbackslash frenchspacing, which will reduce the space after periods. Next, move all your figures and tables to the top of the page. Consider removing less important portions of a figure. If you use \textbackslash centering instead of \textbackslash begin\{center\} in your figure environment, you can also buy some space. For mathematical environments, you may reduce fontsize {\bf but not below 6.5 point}.


Commands that alter page layout are forbidden. These include \textbackslash columnsep,  \textbackslash float, \textbackslash topmargin, \textbackslash topskip, \textbackslash textheight, \textbackslash textwidth, \textbackslash oddsidemargin, and \textbackslash evensizemargin (this list is not exhaustive). If you alter page layout, you will be required to pay the page fee. Other commands that are questionable and may cause your paper to be rejected include \textbackslash parindent, and \textbackslash parskip. Commands that alter the space between sections are forbidden. The title sec package is not allowed. Regardless of the above, if your paper is obviously ``squeezed" it is not going to to be accepted. Options for reducing the length of a paper include reducing the size of your graphics, cutting text, or paying the extra page charge (if it is offered).


\subsection{Type Font and Size}
Your paper must be formatted in Times Roman or Nimbus. We will not accept papers formatted using Computer Modern or Palatino or some other font as the text or heading typeface. Sans serif, when used, should be Courier. Use Symbol or Lucida or Computer Modern for \textit{mathematics only. }

Do not use type 3 fonts for any portion of your paper, including graphics. Type 3 bitmapped fonts are designed for fixed resolution printers. Most print at 300 dpi even if the printer resolution is 1200 dpi or higher. They also often cause high resolution imagesetter devices to crash. Consequently, AAAI will not accept electronic files containing obsolete type 3 fonts. Files containing those fonts (even in graphics) will be rejected. (Authors using blackboard symbols must avoid packages that use type 3 fonts.)

Fortunately, there are effective workarounds that will prevent your file from embedding type 3 bitmapped fonts. The easiest workaround is to use the required times, helvet, and courier packages with \LaTeX{}2e. (Note that papers formatted in this way will still use Computer Modern for the mathematics. To make the math look good, you'll either have to use Symbol or Lucida, or you will need to install type 1 Computer Modern fonts --- for more on these fonts, see the section ``Obtaining Type 1 Computer Modern.")

If you are unsure if your paper contains type 3 fonts, view the PDF in Acrobat Reader. The Properties/Fonts window will display the font name, font type, and encoding properties of all the fonts in the document. If you are unsure if your graphics contain type 3 fonts (and they are PostScript or encapsulated PostScript documents), create PDF versions of them, and consult the properties window in Acrobat Reader.

The default size for your type must be ten-point with twelve-point leading (line spacing). Start all pages (except the first) directly under the top margin. (See the next section for instructions on formatting the title page.) Indent ten points when beginning a new paragraph, unless the paragraph begins directly below a heading or subheading.


\subsubsection{Obtaining Type 1 Computer Modern for \LaTeX{}.}

If you use Computer Modern for the mathematics in your paper (you cannot use it for the text) you may need to download type 1 Computer fonts. They are available without charge from the American Mathematical Society:
http://www.ams.org/tex/type1-fonts.html.

\subsubsection{Nonroman Fonts.}
If your paper includes symbols in other languages (such as, but not limited to, Arabic, Chinese, Hebrew, Japanese, Thai, Russian and other Cyrillic languages), you must restrict their use to bit-mapped figures.

\subsection{Title and Authors}
Your title must appear centered over both text columns in sixteen-point bold type (twenty-four point leading). The title must be written in Title Case according to the Chicago Manual of Style rules. The rules are a bit involved, but in general verbs (including short verbs like be, is, using, and go), nouns, adverbs, adjectives, and pronouns should be capitalized, (including both words in hyphenated terms), while articles, conjunctions, and prepositions are lower case unless they directly follow a colon or long dash. You can use the online tool \url{https://titlecaseconverter.com/} to double-check the proper capitalization (select the "Chicago" style and mark the "Show explanations" checkbox).

Author's names should appear below the title of the paper, centered in twelve-point type (with fifteen point leading), along with affiliation(s) and complete address(es) (including electronic mail address if available) in nine-point roman type (the twelve point leading). You should begin the two-column format when you come to the abstract.

\subsubsection{Formatting Author Information.}
Author information has to be set according to the following specification depending if you have one or more than one affiliation. You may not use a table nor may you employ the \textbackslash authorblk.sty package. For one or several authors from the same institution, please separate them with commas and write all affiliation directly below (one affiliation per line) using the macros \textbackslash author and \textbackslash affiliations:

\begin{quote}\begin{scriptsize}\begin{verbatim}
\author{
    Author 1, ..., Author n\\
}
\affiliations {
    Address line\\
    ... \\
    Address line\\
}
\end{verbatim}\end{scriptsize}\end{quote}


\noindent For authors from different institutions, use \textbackslash textsuperscript \{\textbackslash rm x \} to match authors and affiliations. Notice that there should not be any spaces between the author name (or comma following it) and the superscript.

\begin{quote}\begin{scriptsize}\begin{verbatim}
\author{
    AuthorOne,\equalcontrib\textsuperscript{\rm 1,\rm 2}
    AuthorTwo,\equalcontrib\textsuperscript{\rm 2}
    AuthorThree,\textsuperscript{\rm 3}\\
    AuthorFour,\textsuperscript{\rm 4}
    AuthorFive \textsuperscript{\rm 5}}
}
\affiliations {
    \textsuperscript{\rm 1}AffiliationOne,\\
    \textsuperscript{\rm 2}AffiliationTwo,\\
    \textsuperscript{\rm 3}AffiliationThree,\\
    \textsuperscript{\rm 4}AffiliationFour,\\
    \textsuperscript{\rm 5}AffiliationFive\\
    \{email, email\}@affiliation.com,
    email@affiliation.com,
    email@affiliation.com,
    email@affiliation.com
}
\end{verbatim}\end{scriptsize}\end{quote}

You can indicate that some authors contributed equally using the \textbackslash equalcontrib command. This will add a marker after the author names and a footnote on the first page.

Note that you may want to  break the author list for better visualization. You can achieve this using a simple line break (\textbackslash  \textbackslash).

\subsection{\LaTeX{} Copyright Notice}
The copyright notice automatically appears if you use aaai24.sty. It has been hardcoded and may not be disabled.

\subsection{Credits}
Any credits to a sponsoring agency should appear in the acknowledgments section, unless the agency requires different placement. If it is necessary to include this information on the front page, use
\textbackslash thanks in either the \textbackslash author or \textbackslash title commands.
For example:
\begin{quote}
\begin{small}
\textbackslash title\{Very Important Results in AI\textbackslash thanks\{This work is
 supported by everybody.\}\}
\end{small}
\end{quote}
Multiple \textbackslash thanks commands can be given. Each will result in a separate footnote indication in the author or title with the corresponding text at the botton of the first column of the document. Note that the \textbackslash thanks command is fragile. You will need to use \textbackslash protect.

Please do not include \textbackslash pubnote commands in your document.

\subsection{Abstract}
Follow the example commands in this document for creation of your abstract. The command \textbackslash begin\{abstract\} will automatically indent the text block. Please do not indent it further. {Do not include references in your abstract!}

\subsection{Page Numbers}

Do not print any page numbers on your paper. The use of \textbackslash pagestyle is forbidden.

\subsection{Text}
The main body of the paper must be formatted in black, ten-point Times Roman with twelve-point leading (line spacing). You may not reduce font size or the linespacing. Commands that alter font size or line spacing (including, but not limited to baselinestretch, baselineshift, linespread, and others) are expressly forbidden. In addition, you may not use color in the text.

\subsection{Citations}
Citations within the text should include the author's last name and year, for example (Newell 1980). Append lower-case letters to the year in cases of ambiguity. Multiple authors should be treated as follows: (Feigenbaum and Engelmore 1988) or (Ford, Hayes, and Glymour 1992). In the case of four or more authors, list only the first author, followed by et al. (Ford et al. 1997).

\subsection{Extracts}
Long quotations and extracts should be indented ten points from the left and right margins.

\begin{quote}
This is an example of an extract or quotation. Note the indent on both sides. Quotation marks are not necessary if you offset the text in a block like this, and properly identify and cite the quotation in the text.

\end{quote}

\subsection{Footnotes}
Use footnotes judiciously, taking into account that they interrupt the reading of the text. When required, they should be consecutively numbered throughout with superscript Arabic numbers. Footnotes should appear at the bottom of the page, separated from the text by a blank line space and a thin, half-point rule.

\subsection{Headings and Sections}
When necessary, headings should be used to separate major sections of your paper. Remember, you are writing a short paper, not a lengthy book! An overabundance of headings will tend to make your paper look more like an outline than a paper. The aaai24.sty package will create headings for you. Do not alter their size nor their spacing above or below.

\subsubsection{Section Numbers.}
The use of section numbers in AAAI Press papers is optional. To use section numbers in \LaTeX{}, uncomment the setcounter line in your document preamble and change the 0 to a 1. Section numbers should not be used in short poster papers and/or extended abstracts.

\subsubsection{Section Headings.}
Sections should be arranged and headed as follows:
\begin{enumerate}
\item Main content sections
\item Appendices (optional)
\item Ethical Statement (optional, unnumbered)
\item Acknowledgements (optional, unnumbered)
\item References (unnumbered)
\end{enumerate}

\subsubsection{Appendices.}
Any appendices must appear after the main content. If your main sections are numbered, appendix sections must use letters instead of arabic numerals. In \LaTeX{} you can use the \texttt{\textbackslash appendix} command to achieve this effect and then use \texttt{\textbackslash section\{Heading\}} normally for your appendix sections.

\subsubsection{Ethical Statement.}
You can write a statement about the potential ethical impact of your work, including its broad societal implications, both positive and negative. If included, such statement must be written in an unnumbered section titled \emph{Ethical Statement}.

\subsubsection{Acknowledgments.}
The acknowledgments section, if included, appears right before the references and is headed ``Acknowledgments". It must not be numbered even if other sections are (use \texttt{\textbackslash section*\{Acknowledgements\}} in \LaTeX{}). This section includes acknowledgments of help from associates and colleagues, credits to sponsoring agencies, financial support, and permission to publish. Please acknowledge other contributors, grant support, and so forth, in this section. Do not put acknowledgments in a footnote on the first page. If your grant agency requires acknowledgment of the grant on page 1, limit the footnote to the required statement, and put the remaining acknowledgments at the back. Please try to limit acknowledgments to no more than three sentences.

\subsubsection{References.}
The references section should be labeled ``References" and must appear at the very end of the paper (don't end the paper with references, and then put a figure by itself on the last page). A sample list of references is given later on in these instructions. Please use a consistent format for references. Poorly prepared or sloppy references reflect badly on the quality of your paper and your research. Please prepare complete and accurate citations.

\subsection{Illustrations and  Figures}



% Using the \centering command instead of \begin{center} ... \end{center} will save space
% Positioning your figure at the top of the page will save space and make the paper more readable
% Using 0.95\columnwidth in conjunction with the


Your paper must compile in PDF\LaTeX{}. Consequently, all your figures must be .jpg, .png, or .pdf. You may not use the .gif (the resolution is too low), .ps, or .eps file format for your figures.

Figures, drawings, tables, and photographs should be placed throughout the paper on the page (or the subsequent page) where they are first discussed. Do not group them together at the end of the paper. If placed at the top of the paper, illustrations may run across both columns. Figures must not invade the top, bottom, or side margin areas. Figures must be inserted using the \textbackslash usepackage\{graphicx\}. Number figures sequentially, for example, figure 1, and so on. Do not use minipage to group figures.

If you normally create your figures using pgfplots, please create the figures first, and then import them as pdfs with proper bounding boxes, as the bounding and trim boxes created by pfgplots are fragile and not valid.

When you include your figures, you must crop them \textbf{outside} of \LaTeX{}. The command \textbackslash includegraphics*[clip=true, viewport 0 0 10 10]{...} might result in a PDF that looks great, but the image is \textbf{not really cropped.} The full image can reappear (and obscure whatever it is overlapping) when page numbers are applied or color space is standardized. Figures \ref{fig1}, and \ref{fig2} display some unwanted results that often occur.

If your paper includes illustrations that are not compatible with PDF\TeX{} (such as .eps or .ps documents), you will need to convert them. The epstopdf package will usually work for eps files. You will need to convert your ps files to PDF in either case.

\subsubsection {Figure Captions.}The illustration number and caption must appear \textit{under} the illustration. Labels and other text with the actual illustration must be at least nine-point type. However, the font and size of figure captions must be 10 point roman. Do not make them smaller, bold, or italic. (Individual words may be italicized if the context requires differentiation.)

\subsection{Tables}

Tables should be presented in 10 point roman type. If necessary, they may be altered to 9 point type. You may not use any commands that further reduce point size below nine points. Tables that do not fit in a single column must be placed across double columns. If your table won't fit within the margins even when spanning both columns, you must split it. Do not use minipage to group tables.

\subsubsection {Table Captions.} The number and caption for your table must appear \textit{under} (not above) the table.  Additionally, the font and size of table captions must be 10 point roman and must be placed beneath the figure. Do not make them smaller, bold, or italic. (Individual words may be italicized if the context requires differentiation.)



\subsubsection{Low-Resolution Bitmaps.}
You may not use low-resolution (such as 72 dpi) screen-dumps and GIF files---these files contain so few pixels that they are always blurry, and illegible when printed. If they are color, they will become an indecipherable mess when converted to black and white. This is always the case with gif files, which should never be used. The resolution of screen dumps can be increased by reducing the print size of the original file while retaining the same number of pixels. You can also enlarge files by manipulating them in software such as PhotoShop. Your figures should be 300 dpi when incorporated into your document.

\subsubsection{\LaTeX{} Overflow.}
\LaTeX{} users please beware: \LaTeX{} will sometimes put portions of the figure or table or an equation in the margin. If this happens, you need to make the figure or table span both columns. If absolutely necessary, you may reduce the figure, or reformat the equation, or reconfigure the table.{ \bf Check your log file!} You must fix any overflow into the margin (that means no overfull boxes in \LaTeX{}). \textbf{Nothing is permitted to intrude into the margin or gutter.}


\subsubsection{Using Color.}
Use of color is restricted to figures only. It must be WACG 2.0 compliant. (That is, the contrast ratio must be greater than 4.5:1 no matter the font size.) It must be CMYK, NOT RGB. It may never be used for any portion of the text of your paper. The archival version of your paper will be printed in black and white and grayscale. The web version must be readable by persons with disabilities. Consequently, because conversion to grayscale can cause undesirable effects (red changes to black, yellow can disappear, and so forth), we strongly suggest you avoid placing color figures in your document. If you do include color figures, you must (1) use the CMYK (not RGB) colorspace and (2) be mindful of readers who may happen to have trouble distinguishing colors. Your paper must be decipherable without using color for distinction.

\subsubsection{Drawings.}
We suggest you use computer drawing software (such as Adobe Illustrator or, (if unavoidable), the drawing tools in Microsoft Word) to create your illustrations. Do not use Microsoft Publisher. These illustrations will look best if all line widths are uniform (half- to two-point in size), and you do not create labels over shaded areas. Shading should be 133 lines per inch if possible. Use Times Roman or Helvetica for all figure call-outs. \textbf{Do not use hairline width lines} --- be sure that the stroke width of all lines is at least .5 pt. Zero point lines will print on a laser printer, but will completely disappear on the high-resolution devices used by our printers.

\subsubsection{Photographs and Images.}
Photographs and other images should be in grayscale (color photographs will not reproduce well; for example, red tones will reproduce as black, yellow may turn to white, and so forth) and set to a minimum of 300 dpi. Do not prescreen images.

\subsubsection{Resizing Graphics.}
Resize your graphics \textbf{before} you include them with LaTeX. You may \textbf{not} use trim or clip options as part of your \textbackslash includegraphics command. Resize the media box of your PDF using a graphics program instead.

\subsubsection{Fonts in Your Illustrations.}
You must embed all fonts in your graphics before including them in your LaTeX document.

\subsubsection{Algorithms.}
Algorithms and/or programs are a special kind of figures. Like all illustrations, they should appear floated to the top (preferably) or bottom of the page. However, their caption should appear in the header, left-justified and enclosed between horizontal lines, as shown in Algorithm~\ref{alg:algorithm}. The algorithm body should be terminated with another horizontal line. It is up to the authors to decide whether to show line numbers or not, how to format comments, etc.

In \LaTeX{} algorithms may be typeset using the {\tt algorithm} and {\tt algorithmic} packages, but you can also use one of the many other packages for the task.

\begin{algorithm}[tb]
\caption{Example algorithm}
\label{alg:algorithm}
\textbf{Input}: Your algorithm's input\\
\textbf{Parameter}: Optional list of parameters\\
\textbf{Output}: Your algorithm's output
\begin{algorithmic}[1] %[1] enables line numbers
\STATE Let $t=0$.
\WHILE{condition}
\STATE Do some action.
\IF {conditional}
\STATE Perform task A.
\ELSE
\STATE Perform task B.
\ENDIF
\ENDWHILE
\STATE \textbf{return} solution
\end{algorithmic}
\end{algorithm}

\subsubsection{Listings.}
Listings are much like algorithms and programs. They should also appear floated to the top (preferably) or bottom of the page. Listing captions should appear in the header, left-justified and enclosed between horizontal lines as shown in Listing~\ref{lst:listing}. Terminate the body with another horizontal line and avoid any background color. Line numbers, if included, must appear within the text column.

\begin{listing}[tb]%
\caption{Example listing {\tt quicksort.hs}}%
\label{lst:listing}%
\begin{lstlisting}[language=Haskell]
quicksort :: Ord a => [a] $\rightarrow$ [a]
quicksort []     = []
quicksort (p:xs) = (quicksort lesser) ++ [p] ++ (quicksort greater)
	where
		lesser  = filter (< p) xs
		greater = filter (>= p) xs
\end{lstlisting}
\end{listing}

\subsection{References}
The AAAI style includes a set of definitions for use in formatting references with BibTeX. These definitions make the bibliography style fairly close to the ones  specified in the Reference Examples appendix below. To use these definitions, you also need the BibTeX style file ``aaai24.bst," available in the AAAI Author Kit on the AAAI web site. Then, at the end of your paper but before \textbackslash end{document}, you need to put the following lines:

\begin{quote}
\begin{small}
\textbackslash bibliography\{bibfile1,bibfile2,...\}
\end{small}
\end{quote}

Please note that the aaai24.sty class already sets the bibliographystyle for you, so you do not have to place any \textbackslash bibliographystyle command in the document yourselves. The aaai24.sty file is incompatible with the hyperref and navigator packages. If you use either, your references will be garbled and your paper will be returned to you.

References may be the same size as surrounding text. However, in this section (only), you may reduce the size to \textbackslash small if your paper exceeds the allowable number of pages. Making it any smaller than 9 point with 10 point linespacing, however, is not allowed. A more precise and exact method of reducing the size of your references minimally is by means of the following command: \begin{quote}
\textbackslash fontsize\{9.8pt\}\{10.8pt\}
\textbackslash selectfont\end{quote}

\noindent You must reduce the size equally for both font size and line spacing, and may not reduce the size beyond \{9.0pt\}\{10.0pt\}.

The list of files in the \textbackslash bibliography command should be the names of your BibTeX source files (that is, the .bib files referenced in your paper).

The following commands are available for your use in citing references:
\begin{quote}
{\em \textbackslash cite:} Cites the given reference(s) with a full citation. This appears as ``(Author Year)'' for one reference, or ``(Author Year; Author Year)'' for multiple references.\smallskip\\
{\em \textbackslash shortcite:} Cites the given reference(s) with just the year. This appears as ``(Year)'' for one reference, or ``(Year; Year)'' for multiple references.\smallskip\\
{\em \textbackslash citeauthor:} Cites the given reference(s) with just the author name(s) and no parentheses.\smallskip\\
{\em \textbackslash citeyear:} Cites the given reference(s) with just the date(s) and no parentheses.
\end{quote}
You may also use any of the \emph{natbib} citation commands.


\section{Proofreading Your PDF}
Please check all the pages of your PDF file. The most commonly forgotten element is the acknowledgements --- especially the correct grant number. Authors also commonly forget to add the metadata to the source, use the wrong reference style file, or don't follow the capitalization rules or comma placement for their author-title information properly. A final common problem is text (expecially equations) that runs into the margin. You will need to fix these common errors before submitting your file.

\section{Improperly Formatted Files }
In the past, AAAI has corrected improperly formatted files submitted by the authors. Unfortunately, this has become an increasingly burdensome expense that we can no longer absorb). Consequently, if your file is improperly formatted, it will be returned to you for correction.

\section{Naming Your Electronic File}
We require that you name your \LaTeX{} source file with the last name (family name) of the first author so that it can easily be differentiated from other submissions. Complete file-naming instructions will be provided to you in the submission instructions.

\section{Submitting Your Electronic Files to AAAI}
Instructions on paper submittal will be provided to you in your acceptance letter.

\section{Inquiries}
If you have any questions about the preparation or submission of your paper as instructed in this document, please contact AAAI Press at the address given below. If you have technical questions about implementation of the aaai style file, please contact an expert at your site. We do not provide technical support for \LaTeX{} or any other software package. To avoid problems, please keep your paper simple, and do not incorporate complicated macros and style files.

\begin{quote}
\noindent AAAI Press\\
1900 Embarcadero Road, Suite 101\\
Palo Alto, California 94303-3310 USA\\
\textit{Telephone:} (650) 328-3123\\
\textit{E-mail:} See the submission instructions for your particular conference or event.
\end{quote}

\section{Additional Resources}
\LaTeX{} is a difficult program to master. If you've used that software, and this document didn't help or some items were not explained clearly, we recommend you read Michael Shell's excellent document (testflow doc.txt V1.0a 2002/08/13) about obtaining correct PS/PDF output on \LaTeX{} systems. (It was written for another purpose, but it has general application as well). It is available at www.ctan.org in the tex-archive.

\appendix
\section{Reference Examples}
\label{sec:reference_examples}
\nobibliography*
Formatted bibliographies should look like the following examples. You should use BibTeX to generate the references. Missing fields are unacceptable when compiling references, and usually indicate that you are using the wrong type of entry (BibTeX class).

\paragraph{Book with multiple authors~\nocite{em:86}} Use the \texttt{@book} class.\\[.2em]
\bibentry{em:86}.

\paragraph{Journal and magazine articles~\nocite{r:80, hcr:83}} Use the \texttt{@article} class.\\[.2em]
\bibentry{r:80}.\\[.2em]
\bibentry{hcr:83}.

\paragraph{Proceedings paper published by a society, press or publisher~\nocite{c:83, c:84}} Use the \texttt{@inproceedings} class. You may abbreviate the \emph{booktitle} field, but make sure that the conference edition is clear.\\[.2em]
\bibentry{c:84}.\\[.2em]
\bibentry{c:83}.

\paragraph{University technical report~\nocite{r:86}} Use the \texttt{@techreport} class.\\[.2em]
\bibentry{r:86}.

\paragraph{Dissertation or thesis~\nocite{c:79}} Use the \texttt{@phdthesis} class.\\[.2em]
\bibentry{c:79}.

\paragraph{Forthcoming publication~\nocite{c:21}} Use the \texttt{@misc} class with a \texttt{note="Forthcoming"} annotation.
\begin{quote}
\begin{footnotesize}
\begin{verbatim}
@misc(key,
  [...]
  note="Forthcoming",
)
\end{verbatim}
\end{footnotesize}
\end{quote}
\bibentry{c:21}.

\paragraph{ArXiv paper~\nocite{c:22}} Fetch the BibTeX entry from the "Export Bibtex Citation" link in the arXiv website. Notice it uses the \texttt{@misc} class instead of the \texttt{@article} one, and that it includes the \texttt{eprint} and \texttt{archivePrefix} keys.
\begin{quote}
\begin{footnotesize}
\begin{verbatim}
@misc(key,
  [...]
  eprint="xxxx.yyyy",
  archivePrefix="arXiv",
)
\end{verbatim}
\end{footnotesize}
\end{quote}
\bibentry{c:22}.

\paragraph{Website or online resource~\nocite{c:23}} Use the \texttt{@misc} class. Add the url in the \texttt{howpublished} field and the date of access in the \texttt{note} field:
\begin{quote}
\begin{footnotesize}
\begin{verbatim}
@misc(key,
  [...]
  howpublished="\url{http://...}",
  note="Accessed: YYYY-mm-dd",
)
\end{verbatim}
\end{footnotesize}
\end{quote}
\bibentry{c:23}.

\vspace{.2em}
For the most up to date version of the AAAI reference style, please consult the \textit{AI Magazine} Author Guidelines at \url{https://aaai.org/ojs/index.php/aimagazine/about/submissions#authorGuidelines}

\section{Acknowledgments}
AAAI is especially grateful to Peter Patel Schneider for his work in implementing the original aaai.sty file, liberally using the ideas of other style hackers, including Barbara Beeton. We also acknowledge with thanks the work of George Ferguson for his guide to using the style and BibTeX files --- which has been incorporated into this document --- and Hans Guesgen, who provided several timely modifications, as well as the many others who have, from time to time, sent in suggestions on improvements to the AAAI style. We are especially grateful to Francisco Cruz, Marc Pujol-Gonzalez, and Mico Loretan for the improvements to the Bib\TeX{} and \LaTeX{} files made in 2020.

The preparation of the \LaTeX{} and Bib\TeX{} files that implement these instructions was supported by Schlumberger Palo Alto Research, AT\&T Bell Laboratories, Morgan Kaufmann Publishers, The Live Oak Press, LLC, and AAAI Press. Bibliography style changes were added by Sunil Issar. \verb+\+pubnote was added by J. Scott Penberthy. George Ferguson added support for printing the AAAI copyright slug. Additional changes to aaai24.sty and aaai24.bst have been made by Francisco Cruz and Marc Pujol-Gonzalez.

\bigskip
\noindent Thank you for reading these instructions carefully. We look forward to receiving your electronic files!

\bibliography{aaai24}

\end{document}
